% !TeX root = ../../Book1.tex
% 纲要标题：## 第22章　无限 Galois 理论
% 纲要位置：计划纲要.md，第 767 行。

\chapter{无限 Galois 理论}
\label{b1:ch:22}
\BookChapterNumber{b1:ch:22}

\begin{chapterabstract}
无限代数扩张可由有限子扩张拼合，其自同构须在各有限层上相容。本章把无限 Galois 群表示为有限群的逆极限，建立 Krull 拓扑与紧性，证明中间域同闭子群之间的无限 Galois 对应。
\end{chapterabstract}

\begin{learninggoals}
\begin{itemize}
\item 把无限Galois群写成有限商的逆极限，区分相容族与任意坐标选择。
\item 理解Krull拓扑的基本开集，掌握紧性证明中的有限交性质与选择原理。
\item 运用闭子群对应、开子群判别和正规商群结论，识别删除闭性后的反例。
\item 计算有限域及实数域的绝对Galois群，利用有限层判断连续性和表示的相容性。
\end{itemize}
\end{learninggoals}

依次把四次、八次、十六次单位根添入有理数域，便得到不断增大的分圆域。每一层都有有限 Galois 群；若想在所有层的并上定义一个自同构，在较大层上的取值必须限制成较小层已有的取值。独立任选各层的自同构，通常无法拼成同一个映射。无限 Galois 理论把这种相容性放进群的拓扑中，再问无限扩张的中间域究竟对应哪些子群。有限对应仍在每一层起作用，但直接照搬“任意子群对应中间域”会遗漏闭性条件。

\ProseParagraphBreak

本章的对应定理主要使用有限 Galois 对应、一般代数扩张的嵌入延拓、正规性与可分性，其中有限对应所需内容见\secref{b1:sec:1901}和\secref{b1:sec:1902}。有限域的 Frobenius 计算用于\secref{b1:sec:2204}的典型例子。拓扑术语在使用前定义；任意指标下的紧性证明使用 Zorn 引理。实数域的例子使用第12章引用的\cref{b1:thm:fundamental-theorem-algebra}；其证明位于\secref{b1:sec:E02}，调用实分析的中间值定理、Sylow 理论与有限 Galois 对应。末尾简述连续群上同调与无限 Galois 群的联系。


\ProseParagraphBreak
本章的无限 Galois 扩张可与 Lang 的《Algebra》\cite[Chapter VI, \S 14]{Lang2002}对照。Gille 与 Szamuely 的《Central Simple Algebras and Galois Cohomology》\cite[Sections 4.1--4.2]{GilleSzamuely2017CentralSimpleAlgebras}进一步讨论绝对 Galois 群及连续上同调，使用本章的拓扑群知识，也需要系数模和群上同调。有限群的线性表示与特征标则见 Serre 的《Linear Representations of Finite Groups》\cite[Part I, Chapters 1--2, pp. 3--24]{Serre1977FiniteRepresentations}。该书这部分主要讨论复系数表示；本章的有限域系数表示若经过阶数可被系数域特征整除的有限群，完全可约性和普通迹函数的判别能力就不能照搬复系数情形。

% !TeX root = ../../Book1.tex
% 纲要标题：### 22.1 无限代数扩张
% 纲要位置：计划纲要.md，第 769 行。

\section{无限代数扩张}
\label{b1:sec:2201}

无限代数扩张中的每一个元素，仍落在某个有限子扩张中。困难不在单个元素，而在于不同有限层上的选择必须彼此相容。本节先把自同构写成一族有限层上的自同构，再说明为什么任意选择各层自同构通常不能拼成一个整体。

% 纲要内容（仅作写作计划，尚非教材正文）：
%
% * 有限 Galois 子扩张的有向系统；说明有限层须被全局自同构保持，并区分过渡映射与整体限制映射的满射性
% * 代数闭包的自同构群；解释可分闭包与纯不可分部分，以及自同构群同构为何不能代替扩张的可分性
% * 逐层相容的自同构；证明拼合与所选有限层无关
%

\subsection{有限子扩张的有向系统}
\label{b1:subsec:2201:01}
本章固定代数扩张 \(L/K\)。若它正规且可分，仍称 \(L/K\) 为 Galois 扩张，但不再要求次数有限；这时把它的自同构群记为 \(G=\operatorname{Gal}(L/K)\)。有限层上的群阶仍等于相应的扩张次数，整个无限扩张的域次数与自同构群的普通基数却一般不相等。记录各个有限 Galois 层上的自同构及其限制映射，能够保留有限理论中可以逐层计算的信息；这些有限层还将决定群上的拓扑。

\begin{proposition}\label{b1:prop:finite-galois-layers}
设 \(L/K\) 为代数 Galois 扩张。\(L\) 的每个有限子集都包含在某个有限 Galois 子扩张 \(N/K\) 中。所有这样的 \(N\) 按包含关系构成有向集合，且 \(L=\bigcup_N N\)。
\end{proposition}
\begin{proof}
取有限集 \(a_1,\ldots,a_r\in L\)。每个 \(a_i\) 的最小多项式可分，且由正规性在 \(L\) 中分裂。只有有限多个这样的多项式，每个又只有有限多个根，所以它们的全部根组成有限集。将这些根加入 \(K\)，所得域 \(N\) 仍在 \(L\) 内，由有限多个代数元生成，故 \(N/K\) 有限。把相同的首一最小多项式只保留一次，所得乘积可分，因为不同不可约多项式没有公共根；\(N\) 正是该乘积的分裂域，因此正规且可分，并含原有限集。空集可取 \(N=K\)。

若 \(N_1,N_2\) 是两个有限 Galois 子扩张，各取有限生成元，将它们的并按上一段处理，就有同时包含两域的有限 Galois 子扩张。这就是按包含关系有向的含义；有限次重复，任意有限多个这样的域都能放进一个共同有限 Galois 层中。每个 \(a\in L\) 所成的单点集也被某层包含，故它们的并为 \(L\)。
\end{proof}
有向集合的抽象定义见\cref{b1:def:directed-system}；这里的共同上层已由上述证明构造。记这个有向集合为 \(\mathcal N\)，并记 \(G_N=\operatorname{Gal}(N/K)\)。选取有限层时要求 \(N/K\) 正规，是为了能在它上面记录整体自同构的限制：由\cref{b1:thm:normal-embedding-criterion}，对任意 \(\sigma\in G\)，有 \(\sigma(N)=N\)。若只取任意有限子扩张 \(E/K\)，\(\sigma(E)\) 可能是另一个共轭子域，\(\sigma|_E\) 就不能视为 \(E\) 的自同构。同理，对 \(N\subseteq M\)，\(M\) 的每个 \(K\)-自同构都保持 \(N\)，所以限制给出群同态
\[
\operatorname{res}_{M,N}:G_M\longrightarrow G_N.
\]
这个有限层间的限制映射是否满射，以及 \(G\) 能否实现有限层上的每个自同构，都还需要嵌入延拓定理保证。

\begin{lemma}\label{b1:lem:galois-restriction-surjective}
设 \(L/K\) 为代数 Galois 扩张，\(G=\operatorname{Gal}(L/K)\)。若 \(N\subseteq M\subseteq L\) 都是 \(L/K\) 的有限 Galois 子扩张，并记 \(G_N=\operatorname{Gal}(N/K)\)、\(G_M=\operatorname{Gal}(M/K)\)，则限制同态
\[
G\longrightarrow G_N,
\qquad
\operatorname{res}_{M,N}:G_M\longrightarrow G_N
\]
都满射。
\end{lemma}
\begin{proof}
取包含 \(L\) 的 \(K\) 的代数闭包 \(\Omega\)。给定 \(\tau\in G_N\)，把它看作 \(N\rightarrow\Omega\) 的嵌入。\cref{b1:thm:embedding-extension}应用于代数扩张 \(L/N\)，将它延拓为 \(\sigma:L\rightarrow\Omega\)。因为 \(\sigma|_N=\tau\)，它固定 \(K\)；\cref{b1:thm:normal-embedding-criterion}及 \(L/K\) 的正规性给出 \(\sigma(L)=L\)，故 \(\sigma\in G\)，且它在 \(N\) 上恰为 \(\tau\)。这证明 \(G\rightarrow G_N\) 满射，即有限层的每个自同构都能在整体扩张上实现。

把同一论证中的 \(L\) 换成 \(M\)，\(M/N\) 仍代数而 \(M/K\) 正规，就得到 \(\tau\) 到 \(M\) 的自同构延拓。因此 \(G_M\rightarrow G_N\) 也满射，即较小有限层上的自同构能提升到较大有限 Galois 层。第一种满射控制整体群在每个有限坐标上的像，第二种控制有限层之间的过渡映射。对于三层域，逐次限制等于直接限制。
\end{proof}
这两种满射性来自延拓与正规性，并不是把各有限群写成坐标以后自动得到的性质。

\subsection{代数闭包的自同构群}
\label{b1:subsec:2201:02}
不完美域的代数闭包中，可能有基域之外的元素被每个基域自同构固定。例如 \(K=\mathbb F_p(t)\)，在代数闭包中取 \(u=t^{1/p}\)。多项式 \(X^p-t\) 在 \(K\) 上不可约，且只有 \(u\) 这一个不同根，所以 \(u\notin K\)，每个 \(K\)-自同构却都固定 \(u\)。因此 \(\operatorname{Aut}_K(\overline K)\) 在整个代数闭包上的不动域可以严格大于 \(K\)。要使自同构的作用能够检测扩张中的不同共轭根，需要考察可分部分。

\ProseParagraphBreak
固定 \(K\) 的代数闭包 \(\overline K\)。其中所有在 \(K\) 上可分的元素组成子域，记为 \(K^{\mathrm{sep}}\)，称为 \(K\) 的\term[kefen-bibao]{可分闭包}[separable closure]。它是子域的依据为\cref{b1:prop:separable-subfield}；任一可分元的全部共轭根仍然可分，且已在 \(\overline K\) 中，所以 \(K^{\mathrm{sep}}/K\) 正规而可分。
\symidx{\(K^{\mathrm{sep}}\)：所选代数闭包中的可分闭包}

\begin{definition}\label{b1:def:absolute-galois-group}
设 \(K\) 为域，固定一个代数闭包 \(\overline K\)，并令 \(K^{\mathrm{sep}}\subseteq\overline K\) 为所有在 \(K\) 上可分的元素组成的可分闭包。群
\[
G_K\coloneqq\operatorname{Gal}(K^{\mathrm{sep}}/K)
\]
称为域 \(K\) 的\term[juedui-galois-qun]{绝对 Galois 群}[absolute Galois group]。换一份可分闭包会给出同构的群，但同构一般依赖于闭包之间同构的选择。
\end{definition}

\symidx{\(G_K\)：域的绝对 Galois 群}

\begin{proposition}\label{b1:prop:closure-automorphisms}
设 \(K\) 为域，\(\overline K\) 为 \(K\) 的代数闭包，\(K^{\mathrm{sep}}\subseteq\overline K\) 为其中所有在 \(K\) 上可分的元素组成的子域，并令 \(G_K=\operatorname{Gal}(K^{\mathrm{sep}}/K)\)。则限制映射 \(\operatorname{Aut}_K(\overline K)\rightarrow G_K\) 为群同构。若 \(K\) 完美，则 \(K^{\mathrm{sep}}=\overline K\)；一般情形中不能把 \(\overline K/K\) 直接视为 Galois 扩张。
\end{proposition}
\begin{proof}
\(K\)-自同构保持最小多项式及其可分性，故保持可分闭包，限制映射有定义。若特征为零，所有代数元可分，限制映射就是恒等识别。以下设特征为 \(p>0\)。对任意 \(a\in\overline K\)，\cref{b1:prop:inseparable-polynomial-shape}把它的最小多项式写为
\[
m_{a,K}(X)=g(X^{p^r}),
\]
其中 \(g\) 首一、不可约且可分。代入 \(a\) 得 \(g(a^{p^r})=0\)，故 \(g\) 正是 \(a^{p^r}\) 的最小多项式。因此 \(a^{p^r}\) 在 \(K\) 上可分，属于 \(K^{\mathrm{sep}}\)。这说明每个 \(a\) 都在 \(K^{\mathrm{sep}}\) 上纯不可分。

若两个 \(K\)-自同构 \(\sigma,\tau\) 在 \(K^{\mathrm{sep}}\) 上有相同限制 \(\nu\)，对给定 \(a\) 取上述的 \(r\)，则
\[
\sigma(a)^{p^r}=\nu(a^{p^r})=\tau(a)^{p^r}.
\]
两像都是方程 \(X^{p^r}-\nu(a^{p^r})=0\) 的根。这个方程至多有一个不同根，因为两根之差的 \(p^r\) 次幂为零，而域没有非零幂零元。因此 \(\sigma(a)=\tau(a)\)；对每个 \(a\) 都如此，限制映射便单射。

反过来，给定 \(\nu\in G_K\)，把它看作 \(K^{\mathrm{sep}}\rightarrow\overline K\) 的嵌入。\cref{b1:thm:embedding-extension}应用于代数扩张 \(\overline K/K^{\mathrm{sep}}\)，将它延拓为 \(\sigma:\overline K\rightarrow\overline K\)。这个嵌入固定 \(K\)，代数闭包在 \(K\) 上正规，\cref{b1:thm:normal-embedding-criterion}保证 \(\sigma(\overline K)=\overline K\)，所以它是自同构，得到满射性。完美域的所有代数元可分，最后的等式由定义得到。
\end{proof}
纯不可分部分上的延拓唯一，因而没有增加自同构的选择，这正是两个群同构的原因。但这个同构不会使原本不可分的元素变得可分；这些不能被移动的元素仍在代数闭包中，使该群在那里作用的不动域可以严格大于 \(K\)。所以，还须区分群本身与它所作用的域，不能由上述群同构把 \(\overline K/K\) 直接当作 Galois 扩张。

\subsection{相容自同构}
\label{b1:subsec:2201:03}
从有限层上的自同构恢复整体映射，自然要对 \(a\in L\) 规定 \(\sigma(a)=\sigma_N(a)\)，其中 \(N\) 是任意包含 \(a\) 的有限 Galois 层。首先须解决的是选择无关性：同一个元素可能属于许多层，只有这些层给出的像一致，才能在 \(L\) 上定义一个映射。

\begin{proposition}\label{b1:prop:galois-compatible-family}
设 \(L/K\) 为代数 Galois 扩张，\(G=\operatorname{Gal}(L/K)\)，并令 \(\mathcal N\) 为 \(L/K\) 的全部有限 Galois 子扩张组成的有向集合。对每个 \(N\in\mathcal N\)，记 \(G_N=\operatorname{Gal}(N/K)\)，并给定 \(\sigma_N\in G_N\)。则存在唯一 \(\sigma\in G\) 在每个 \(N\) 上等于 \(\sigma_N\)，当且仅当
\begin{equation}\label{b1:eq:galois-compatibility}
\sigma_M|_N=\sigma_N\qquad(N\subseteq M).
\end{equation}
\end{proposition}
\begin{proof}
限制一个自同构当然满足相容条件。反过来，对 \(a\in L\)，选含 \(a\) 的有限 Galois 层 \(N\)，规定 \(\sigma(a)=\sigma_N(a)\)。若还选了 \(M\)，由\cref{b1:prop:finite-galois-layers}取同时包含两者的一层 \(P\)，相容性给出 \(\sigma_N(a)=\sigma_P(a)=\sigma_M(a)\)，所以定义无关选择。

对任意两个元素，取包含它们的同一有限层，就可在该层验证保持和、积与单位，以及固定 \(K\)。还须保证所得映射可逆。若 \(N\subseteq M\)，相容性给出 \(\sigma_M|_N=\sigma_N\)，而 \(\sigma_N\) 是 \(N\) 的自同构；所以 \(\sigma_M^{-1}\) 在 \(N\) 上的限制正是 \(\sigma_N^{-1}\)。因此逆族 \((\sigma_N^{-1})_N\) 也相容，同样拼成一个映射，在每层上都是 \(\sigma\) 的双边逆，因而在 \(L\) 上也是双边逆。所得 \(\sigma\) 是自同构；\(L\) 是各层的并，全部限制也就唯一确定了它。
\end{proof}
在每层能够延拓，不意味着可以把所有层上任意选定的自同构拼在一起。例如在有限域的塔 \(\mathbb F_q\subset\mathbb F_{q^2}\subset\mathbb F_{q^4}\) 中，取四次层上的 Frobenius \(F(x)=x^q\)。它在二次层上的限制仍为 \(x\mapsto x^q\)；由\cref{b1:thm:finite-field-automorphisms}，这个限制的阶为 \(2\)，是二次扩张的非平凡自同构。所以若预先在二次层选择恒等，它就不能与四次层上选定的 \(F\) 相容，违反\eqref{b1:eq:galois-compatibility}。各层的选择分别合法，还须满足限制相容条件才能拼成整体。下一节把相容族组成的群称为逆极限，并在其上建立适合这些有限信息的拓扑。

% !TeX root = ../../Book1.tex
% 纲要标题：### 22.2 Krull 拓扑
% 纲要位置：计划纲要.md，第 775 行。

\section{Krull 拓扑}
\label{b1:sec:2202}

一个自同构在有限层上的取值，只提供有限的信息。把“在某个有限层上相同”作为接近的标准，就得到 Krull 拓扑。它不是额外附加的距离，而是把有限 Galois 群之间的相容关系变成可用的开集。所需拓扑概念及紧性证明均在本节给出。

% 纲要内容（仅作写作计划，尚非教材正文）：
%
% * 有限商与逆极限；区分一般有向系统与可数序列，坐标投影的满射性须另行证明
% * Galois 群的拓扑；把基本开陪集解释为在某一有限层上作用相同，并证明有限元素观察给出同一拓扑
% * 就地介绍离散拓扑、积拓扑、有限交性质及闭子空间，完整证明有限离散空间积的紧性；展开生成滤子、极大滤子的二择一性质与有限划分取值，并说明选择原理的位置
% * 把逆极限识别为相容条件定义的闭子空间，连续证明 Galois 群的紧性、全不连通性与开子群性质
% * 附录 A 仅汇总和补充这段拓扑，不承担本章缺失的必需引理
%

\subsection{有限商与逆极限}
\label{b1:subsec:2202:01}
有限 Galois 层按域包含关系排列：\(N\subseteq M\) 时，\(M\) 记录更多元素上的取值，而限制映射却从较大层的群 \(G_M\) 指向较小层的群 \(G_N\)。下面的逆系统沿用这个方向，指标向后增大，坐标信息则限制到前层。

\begin{definition}\label{b1:def:inverse-limit}
设 \((I,\leq)\) 为有向集。对每个 \(i\in I\) 给定群 \(G_i\)，并对 \(i\leq j\) 给定同态 \(r_{ji}:G_j\rightarrow G_i\)，满足 \(r_{ii}=\operatorname{id}_{G_i}\)，且对所有 \(i\leq j\leq k\) 都有 \(r_{ki}=r_{ji}r_{kj}\)。在直积 \(\prod_iG_i\) 中取相容子群
\[
\varprojlim_iG_i
=\Bigl\{\,(g_i)_i\in\prod_iG_i\mid
r_{ji}(g_j)=g_i\text{ 每当 }i\leq j\,\Bigr\},
\]
并逐坐标定义乘法。称这个群为所给系统的\term[nijixian]{逆极限}[inverse limit]。
\end{definition}
\symidx{\(\varprojlim_iG_i\)：群的逆极限}
同态保持乘法及逆元，故相容族的乘积与逆仍相容；各坐标单位组成单位族。因此这确实是直积的子群。投影到第 \(i\) 个坐标记为 \(\pi_i\)。逆极限把相互一致的各层信息看作一个整体：若群 \(H\) 上给出同态 \(f_i:H\rightarrow G_i\)，且 \(r_{ji}f_j=f_i\) 对每个 \(i\le j\) 成立，则映射 \(h\mapsto\bigl(f_i(h)\bigr)_i\) 的像落在逆极限中。它保持群运算，并且是唯一满足 \(\pi_i f=f_i\) 的同态 \(f:H\rightarrow\varprojlim_iG_i\)，因为每个坐标已经指定。在本章取 \(H=G\)、\(f_N\) 为限制映射，正是把整体自同构的有限层信息放进逆极限。其底集合构造与唯一性可对照\cref{b1:prop:inverse-limit-set-universal}。

\ProseParagraphBreak
例如，通常的模数约化给出逆系统
\[
\mathbb Z/2\mathbb Z\longleftarrow
\mathbb Z/4\mathbb Z\longleftarrow
\mathbb Z/8\mathbb Z\longleftarrow\cdots.
\]
模 \(2^r\) 的剩余类 \([a]_{2^r}\) 在下一层有两个提升：\([a]_{2^{r+1}}\) 与 \([a+2^r]_{2^{r+1}}\)。逐层相容地选定剩余类，就得到逆极限 \(\mathbb Z_2=\varprojlim_r\mathbb Z/2^r\mathbb Z\) 的一个元素。这个例子是一条可数链；一般有向逆系统未必能化为这样的可数塔，也未必存在一个“下一层”。在 Galois 系统中，任意两个已经考虑的有限层都能进入某个共同上层，各个限制必须在较小层上一致；重复取共同上层，可以同时比较有限多个坐标。

\ProseParagraphBreak
对上一节的有限 Galois 层应用\cref{b1:prop:galois-compatible-family}，得到群同构
\begin{equation}\label{b1:eq:galois-inverse-limit}
G\longrightarrow\varprojlim_{N\in\mathcal N}G_N,
\qquad\sigma\longmapsto(\sigma|_N)_N.
\end{equation}
每个坐标投影满射的依据是\cref{b1:lem:galois-restriction-surjective}：给定 \(G_N\) 中的自同构，嵌入延拓与 \(L/K\) 的正规性使它延拓为 \(G\) 中的自同构，因而在逆极限中得到一个具有指定第 \(N\) 坐标的相容族。一般逆极限的定义并不保证坐标投影满射；它只收集能够扩成完整相容族的坐标，没有保证某一层的每个取值都能如此扩充。今后把两侧识别，不仅记录抽象群运算，还保留全部限制映射。

\subsection{Galois 群的拓扑}
\label{b1:subsec:2202:02}
设 \(X\) 为集合，\(\mathcal T\) 是 \(X\) 的一族子集。若 \(\mathcal T\) 包含 \(\varnothing,X\)，并且对任意并和有限交封闭，则称 \(\mathcal T\) 为 \(X\) 上的一个\term[tuopu]{拓扑}[topology]；\(\mathcal T\) 中的集合称为开集，其补集称为闭集。若 \(\mathcal T=\mathcal P(X)\)，即 \(X\) 的每个子集都开，则称这个拓扑为\term[lisan-tuopu]{离散拓扑}[discrete topology]。设 \(X,Y\) 为拓扑空间，\(f:X\rightarrow Y\) 为映射。若 \(Y\) 中每个开集的原像都是 \(X\) 中的开集，则称 \(f\) 为连续映射；若 \(f\) 还是双射，并且 \(f\) 与 \(f^{-1}\) 都连续，则称 \(f\) 为同胚。本节的每个有限群 \(G_N\) 都采用离散拓扑。

\ProseParagraphBreak
设 \((X_i)_{i\in I}\) 为拓扑空间族，\(P=\prod_{i\in I}X_i\)，并以 \(p_i:P\rightarrow X_i\) 表示坐标投影。由所有集合
\[
\bigcap_{i\in F}p_i^{-1}(U_i),
\qquad F\subseteq I\text{ 有限},\quad U_i\subseteq X_i\text{ 开},
\]
的任意并组成的拓扑称为 \(P\) 上的\term[ji-tuopu]{积拓扑}[product topology]，其中 \(F=\varnothing\) 时交集约定为 \(P\)；所列有限交称为基本开集。对任意子集 \(Y\subseteq P\)，若把 \(Y\cap U\)，其中 \(U\) 遍历 \(P\) 的开集，规定为 \(Y\) 中的开集，则称所得拓扑为 \(Y\) 的子空间拓扑。应用于有限离散群族 \((G_N)\) 时，基本开集只规定有限多个坐标的取值，其余坐标任意。

\begin{samepage}
\begin{definition}\label{b1:def:krull-topology}
设 \(L/K\) 为代数 Galois 扩张，\(G=\operatorname{Gal}(L/K)\)，并令 \(\mathcal N\) 为 \(L/K\) 的全部有限 Galois 子扩张组成的有向集合。对 \(N\in\mathcal N\)，记 \(G_N=\operatorname{Gal}(N/K)\)，并令 \(\pi_N:G\rightarrow G_N\) 为限制映射。通过同构 \(G\cong\varprojlim_{N\in\mathcal N}G_N\)，把 \(G\) 看作有限离散群直积 \(\prod_NG_N\) 的子空间。由该子空间拓扑得到的 \(G\) 上的拓扑称为\term[krull-tuopu]{Krull 拓扑}[Krull topology]。若
\[
U_N\coloneqq\operatorname{Gal}(L/N)=\ker\pi_N,
\]
则对每个 \(\sigma\in G\)，其基本开邻域为 \(\sigma U_N\)，其中 \(N\) 遍历有限 Galois 子扩张。
\end{definition}
\end{samepage}
\symidx{\(U_N=\operatorname{Gal}(L/N)\)：Krull 拓扑的基本开正规子群}
对 \(\sigma,\tau\in G\)，陪集的具体含义是
\[
\tau\in\sigma U_N
\iff\pi_N(\tau)=\pi_N(\sigma)
\iff\tau|_N=\sigma|_N.
\]
也就是说，在有限层 \(N\) 上无法区分 \(\tau\) 与 \(\sigma\)。一个基本开集原本可能限制有限多个坐标 \(N_1,\ldots,N_r\)，但\cref{b1:prop:finite-galois-layers}给出同时包含这些域的有限 Galois 层 \(M\)。于是 \(\sigma U_M\subseteq\bigcap_j\sigma U_{N_j}\)，只规定第 \(M\) 个坐标就得到包含 \(\sigma\) 且落在原基本开集内的邻域。因此所列陪集确实构成邻域基。它们也都是闭集：其补集是同一有限商中其余陪集的并。

\begin{proposition}\label{b1:prop:krull-group-operations}
设 \(L/K\) 为代数 Galois 扩张，\(G=\operatorname{Gal}(L/K)\)，并在 \(G\) 上赋予 Krull 拓扑。则乘法 \(G\times G\rightarrow G\) 与取逆 \(G\rightarrow G\) 都连续，所以 \(G\) 为拓扑群。
\end{proposition}
\begin{proof}
若要确定 \(\sigma\tau\) 在 \(N\) 上的取值，只需分别确定 \(\sigma,\tau\) 在 \(N\) 上的取值，因为限制保持乘法。具体地，\(\sigma'\in\sigma U_N\)、\(\tau'\in\tau U_N\) 使 \(\pi_N(\sigma'\tau')=\pi_N(\sigma\tau)\)，故它们的乘积落在 \(\sigma\tau U_N\) 中。这在每个点给出乘法连续性。限制也保持逆元，因此确定 \(\sigma|_N\) 就确定 \(\sigma^{-1}|_N\)，得到取逆连续。
\end{proof}
于是左右平移均为同胚。若只要求 \(\tau(a_j)=\sigma(a_j)\) 对有限多个 \(a_1,\ldots,a_r\in L\) 成立，取包含它们的有限 Galois 层 \(N\)，则 \(\sigma U_N\) 包含于这个有限观察集合；该集合的每个点都有这样的陪集邻域，所以它开。反过来，取 \(N/K\) 的一组有限基 \(b_1,\ldots,b_s\)。自同构固定 \(K\) 并保持线性组合，故在这些基元素上取值相同就等价于在整个 \(N\) 上取值相同，所得集合正是 \(\sigma U_N\)。因此“在有限多个元素上取值相同”与“在某个有限 Galois 层上取值相同”给出相同拓扑；前者表示有限观察，后者便于使用有限群的结构。

\subsection{有限离散空间积的紧性}
\label{b1:subsec:2202:04}
相容数据的条件可能有无限多个，我们只能一次检查有限批条件。对于闭条件，紧性将把“每一有限批条件都有共同解”加强为“全部条件有共同解”，这也是后面处理闭子群与商群拓扑时需要的性质。在有限离散空间的任意直积中，各坐标只有有限个取值，这将使上述性质能够由有限划分的选择得到。

\ProseParagraphBreak
设 \(X\) 为拓扑空间。若 \(X\) 的每个开覆盖都有有限子覆盖，则称 \(X\) 为\term[nijin-kongjian]{拟紧空间}[quasi-compact space]。若 \(X\) 中任意两个不同的点都有互不相交的开邻域，则称 \(X\) 具有 Hausdorff 性\footnote{豪斯多夫（Felix Hausdorff，1868-1942），德国数学家，研究集合论与拓扑学，建立拓扑空间的系统理论。}；若 \(X\) 同时拟紧并具有 Hausdorff 性，则称 \(X\) 为\term[jin-kongjian]{紧空间}[compact space]。

\ProseParagraphBreak
设 \(X\) 为集合，\(\mathcal C\subseteq\mathcal P(X)\) 为一族子集。若 \(\mathcal C\) 中任意有限多个集合的交都非空，其中空族的交约定为 \(X\)，则称 \(\mathcal C\) 具有\term[youxian-jiao-xingzhi]{有限交性质}[finite intersection property]。对坐标条件而言，每个集合收集满足某一条件的点，例如满足 \(r_{ji}(g_j)=g_i\) 或 \(g_i=a_i\) 的点；有限交性质就是任意有限批要求都能同时满足。

\ProseParagraphBreak
有限子覆盖性质等价于每个具有有限交性质的闭集族都有非空总交。事实上，若这种闭集族的总交为空，其开补集就覆盖 \(X\)；有限子覆盖会使相应闭集的有限交为空，矛盾。反过来，若一个开覆盖没有有限子覆盖，则开集的闭补集具有有限交性质，总交却为空，也与闭集族的条件矛盾。在 Hausdorff 空间中，这个闭集判据便等价于紧性。有限交判据及闭子空间的结论另在\cref{b1:prop:compact-fip}中统一陈述并证明。

\ProseParagraphBreak
设 \(X\) 为集合。若 \(\mathcal F\subseteq\mathcal P(X)\) 包含 \(X\) 而不包含空集，对有限交封闭，并且包含其每个成员的所有超集，则称 \(\mathcal F\) 为\term[zhen-lvzi]{真滤子}[proper filter]。若真滤子在包含关系下极大，则称为\term[chaolvzi]{超滤子}[ultrafilter]。这里的极大是指不能再严格扩大而仍保持为真滤子，并不是包含所有其他真滤子的最大者；尤其不能把空集加入其中。

\ProseParagraphBreak
每个有限离散空间本身都是紧空间，因为有限多个点就能从任意开覆盖中选出有限子覆盖，单点开集又可分离不同点。对任意指标的乘积，我们将把有限条件组成的滤子扩大为超滤子，由它在每个有限坐标中确定一个值，再验证所得点满足全部原条件。

\begin{lemma}\label{b1:lem:finite-discrete-product-compact}
任意一族有限离散空间的直积，在积拓扑下是紧空间。
\end{lemma}
\begin{proof}
记直积为 \(X=\prod_{i\in I}X_i\)。两个不同点必有某个坐标不同，规定该坐标值的两个基本开集将它们分离，故 \(X\) 为 Hausdorff 空间。若 \(X\) 为空，紧性直接成立。以下给定具有有限交性质的闭集族 \((C_\alpha)_\alpha\)，证明总交非空。

把原闭集族的有限交及其超集组成的族记为
\[
\mathcal F=
\left\{\,A\subseteq X\,\middle|\,
C_{\alpha_1}\cap\cdots\cap C_{\alpha_r}\subseteq A
\text{ 对某个有限指标列表成立}\,\right\},
\]
其中允许空列表，其交为 \(X\)。它包含 \(X\)，对取超集封闭；两个成员各包含一个有限交，它们的交便包含这两批闭集的共同有限交，故仍属于 \(\mathcal F\)。有限交性质又保证这些交都非空，所以 \(\varnothing\notin\mathcal F\)。这证明 \(\mathcal F\) 为真滤子，并包含每个 \(C_\alpha\)。

对包含 \(\mathcal F\) 的全体真滤子按包含关系应用\cref{b1:thm:A-zorn}。这些滤子组成 \(\mathcal P(\mathcal P(X))\) 的子集，因而是一个集合；它非空，因为 \(\mathcal F\) 本身在其中。若一条非空链中的滤子取并，该并仍包含 \(X\)，不包含空集，并且对超集封闭。有限多个属于这个并的集合，分别属于链中的有限多个滤子；链全序，故其中有一个滤子同时包含它们，有限交仍属于该滤子，也属于链的并。因此链的并是一个包含 \(\mathcal F\) 的真滤子，作为上界；空链则以 \(\mathcal F\) 为上界。Zorn 引理给出一个极大者 \(\mathcal U\)。任何严格包含它的真滤子也会包含 \(\mathcal F\)，所以这里的极大者确为超滤子。

对任何 \(A\subseteq X\)，\(A\) 与 \(X\setminus A\) 恰有一个属于 \(\mathcal U\)。它们不能同时属于，因为交为空。现设 \(A\notin\mathcal U\)。若有某个 \(B\in\mathcal U\) 与 \(A\) 不相交，则 \(B\subseteq X\setminus A\)，向上封闭性给出 \(X\setminus A\in\mathcal U\)。若反而每个 \(B\in\mathcal U\) 都与 \(A\) 相交，则加入 \(A\) 后的任意有限交为 \(B_1\cap\cdots\cap B_r\cap A\)；前 \(r\) 个集合的交仍属于 \(\mathcal U\)，按假设与 \(A\) 相交，所以这个有限交非空。更具体地，族
\[
\mathcal V=\{\,D\subseteq X\mid B\cap A\subseteq D
\text{ 对某个 }B\in\mathcal U\text{ 成立}\,\}
\]
包含 \(\mathcal U\) 和 \(A\)，对有限交与取超集封闭，且不包含空集，因而是严格包含 \(\mathcal U\) 的真滤子，违背极大性。这就证明了二择一性质。

对于 \(X\) 的任意有限划分，恰有一块属于 \(\mathcal U\)。不同块的交为空，故至多一块属于；若一块也没有，二择一性质使全部块的补集都属于 \(\mathcal U\)，有限求交却为空，矛盾。现在对每个坐标 \(i\)，纤维 \(p_i^{-1}(\{t\})\)、\(t\in X_i\)，恰好给出一个有限划分，因而唯一确定其中属于 \(\mathcal U\) 的一块。将相应的坐标值记为 \(a_i\)，得到点 \(a=(a_i)_i\in X\)。对于每个有限集 \(F\subseteq I\)，柱状集合
\[
B_F(a)=\bigcap_{i\in F}p_i^{-1}(\{a_i\})
\]
属于 \(\mathcal U\)；这些集合构成 \(a\) 的基本开邻域，因为每个 \(X_i\) 都离散。

若 \(a\notin C_\alpha\)，其开补集包含 \(a\)。积拓扑的定义保证它包含一个只限制有限坐标的基本开邻域；各坐标空间离散，再把所限制的坐标收紧到单点，就有某个有限集 \(F\subseteq I\) 满足 \(a\in B_F(a)\subseteq X\setminus C_\alpha\)。可是 \(B_F(a),C_\alpha\) 都属于 \(\mathcal U\)，它们的交却为空，矛盾。所以 \(a\) 属于全部 \(C_\alpha\)，证明总交非空。
\end{proof}
这里借助超滤子同时作出与有限交条件相容的坐标选择。闭子空间的紧性也随即得到：闭子空间中的闭集仍是整个空间中的闭集，有限交性质可以原样用于整个紧空间，分离邻域与子空间相交又保证 Hausdorff 性。

\ProseParagraphBreak
以后还会用到两个基本事实。连续像保留有限子覆盖性质：将像上的开覆盖拉回，取有限子覆盖再映回即可；目标为Hausdorff空间时，像因而紧。Hausdorff 空间中的紧子集闭：给紧子集外一点 \(x\)，对每个子集内点 \(y\) 取分离 \(x,y\) 的开邻域，覆盖紧子集的那些邻域可选有限个；对应的 \(x\) 邻域有限求交，就与整个紧子集不交。因而从紧空间到 Hausdorff 空间的连续双射为同胚：闭集的像紧而闭，使逆映射连续。这些事实汇总为\cref{b1:prop:compact-maps}，将在比较 Galois 商群的拓扑时使用。

\ProseParagraphBreak
上述紧性证明中，选择原理明确用在把有限相容条件生成的真滤子扩充为超滤子的 Zorn 步骤。极大性使它对每个子集与其补集作出二择一决定，进而在每个有限划分中决定唯一一块。每个坐标的值一旦由 \(\mathcal U\) 决定，就是唯一的，不再另作任意选择。代数闭包与嵌入延拓中的 Zorn 引理用于取得极大的部分代数扩张或部分嵌入；这里扩大的是记录相容条件的集合族，两处取得的极大对象不同。

\subsection{紧性、全不连通性与开子群}
\label{b1:subsec:2202:03}
设 \(X\) 为拓扑空间，\(A\subseteq X\) 采用子空间拓扑。若 \(A\) 不能写成两个不相交非空开集的并，则称 \(A\) 为连通子空间。若 \(X\) 的每个非空连通子空间都只有一个点，则称 \(X\) \term[quan-buliantong]{全不连通}[totally disconnected]。

\begin{theorem}\label{b1:thm:krull-compact}
设 \(L/K\) 为代数 Galois 扩张，\(G=\operatorname{Gal}(L/K)\)，并在 \(G\) 上赋予 Krull 拓扑。对 \(L/K\) 的每个有限 Galois 子扩张 \(N/K\)，记 \(U_N=\operatorname{Gal}(L/N)\)。则 \(G\) 是紧、Hausdorff、全不连通的拓扑群。子群 \(H\leq G\) 开，当且仅当它包含某个 \(U_N\)；此时 \(H\) 同时闭且指数有限。
\end{theorem}
\begin{proof}
在有限群直积 \(X=\prod_NG_N\) 中，对 \(N\subseteq M\) 考察映射
\[
\Phi_{M,N}:X\longrightarrow G_N\times G_N,\qquad
(g_P)_P\longmapsto\bigl(\operatorname{res}_{M,N}(g_M),g_N\bigr).
\]
它只依赖两个坐标，目标采用有限离散拓扑，所以连续。相容条件 \(\operatorname{res}_{M,N}(g_M)=g_N\) 定义的集合正是对角集 \(\Delta_N=\{\,(h,h)\mid h\in G_N\,\}\) 的原像；有限离散空间中的对角集闭，故每个相容条件定义闭集。逆极限是这些闭集的交，因此为闭子空间。\cref{b1:lem:finite-discrete-product-compact}及闭子空间的紧性给出 \(G\) 紧。

若 \(\sigma\ne\tau\)，存在 \(a\in L\) 使 \(\sigma(a)\ne\tau(a)\)。由\cref{b1:prop:finite-galois-layers}，\(L\) 是有限 Galois 层之并，可取包含 \(a\) 的有限层 \(N\)，于是 \(\pi_N(\sigma)\ne\pi_N(\tau)\)。两个陪集 \(\sigma U_N,\tau U_N\) 不相交，故 \(G\) 为 Hausdorff 空间。又因 \(\sigma U_N\) 同时开闭，对任何包含 \(\sigma,\tau\) 的子空间 \(A\)，两部分 \(A\cap\sigma U_N\) 与 \(A\setminus\sigma U_N\) 都是非空相对开集，并且不交、合起来为 \(A\)。因此连通子空间不能含两个不同点，得到全不连通性。

若 \(H\) 是开子群，它作为单位元的邻域包含某个 \(U_N\)。反过来，若包含 \(U_N\)，则 \(H\) 为若干个 \(U_N\) 陪集的并，所以开。此时 \(H\) 的指数不超过有限数 \([G:U_N]=|G_N|\)。其补集是其他 \(H\) 陪集的并，而平移保持开集，所以 \(H\) 同时闭。拓扑群运算已由\cref{b1:prop:krull-group-operations}验证。
\end{proof}

全不连通只禁止含有两个不同点的连通子空间，并不要求每个单点都是开集。特别地，若这里的 \(G\) 无限，就不可能离散：否则全体单点开集覆盖 \(G\)，却没有有限子覆盖，与紧性矛盾。因此有限层陪集能够分离不同点，不表示每个点都已经孤立。

\ProseParagraphBreak
设 \(H\) 为拓扑群。若 \(H\) 与某个有限离散群逆系统的逆极限拓扑同构，则称 \(H\) 为\term[toushe-youxian-qun]{投射有限群}[profinite group]；这里拓扑同构指同时为群同构和同胚。\eqref{b1:eq:galois-inverse-limit}说明 \(G\) 就是这样的群。“投射有限”并不意味着整个群只有有限多个元素；有限的是各层群，每个实际出现的坐标像都是 \(H\) 的有限商群。拓扑也是定义的一部分，不能只保留抽象群同构而忘记开邻域；连续性与稠密性都由这份拓扑决定。

% !TeX root = ../../Book1.tex
% 纲要标题：### 22.3 拓扑 Galois 对应
% 纲要位置：计划纲要.md，第 748 行。

\section{拓扑 Galois 对应}
\label{b1:sec:2203}

有限 Galois 对应将全部子群与全部中间域配对。无限情形中，固定一个元素只是在某个有限层上的条件，因此不动域只能区分子群的有限层信息。正确的对应要在群的一侧取闭子群；这个限制既有直接的证明，也有有限域给出的明确反例。

% 纲要内容（仅作写作计划，尚非教材正文）：
%
% * 中间域与闭子群的对应
% * 闭性不可删除的例子
% * 有限子扩张、开子群与正规子群
%

\subsection{中间域与闭子群的对应}
\label{b1:subsec:2203:01}
对中间域 \(K\subseteq E\subseteq L\)，记 \(G_E=\operatorname{Gal}(L/E)\)。对任意子群 \(H\leq G\)，记
\[
L^H=\{\,a\in L\mid\sigma(a)=a\text{ 对每个 }\sigma\in H\,\}.
\]
不动元素对加减乘及非零取逆封闭，故构成中间域。域扩大时固定它的自同构减少，子群扩大时不动域缩小。

\begin{theorem}[无限 Galois 对应]\label{b1:thm:infinite-galois-correspondence}
设 \(L/K\) 为代数 Galois 扩张，\(G=\operatorname{Gal}(L/K)\)，并在 \(G\) 上赋予 Krull 拓扑。映射
\[
E\longmapsto\operatorname{Gal}(L/E),\qquad
H\longmapsto L^H
\]
给出全部中间域与 \(G\) 的全部闭子群之间的反向一一对应。
\end{theorem}
\begin{proof}
先证明 \(G_E\) 闭。对每个 \(a\in E\)，选含 \(a\) 的有限 Galois 层 \(N\)。固定 \(a\) 的自同构，是有限群 \(G_N\) 中固定该元素的那些自同构的原像，因此是有限个 \(U_N\) 陪集的并，既开又闭。\(G_E\) 是对全部 \(a\in E\) 取这些闭子集的交，故闭。

接着证明 \(L^{G_E}=E\)。包含 \(E\subseteq L^{G_E}\) 由定义成立。若 \(a\in L\setminus E\)，则 \(a\) 在 \(E\) 上的最小多项式次数大于一。它整除 \(a\) 在 \(K\) 上的最小多项式，故可分，并由 \(L/K\) 的正规性在 \(L\) 中分裂。取另一个根 \(b\ne a\)，代入最小多项式给出 \(E(a)\rightarrow L\) 的 \(E\)-嵌入 \(a\mapsto b\)。由\cref{b1:thm:embedding-extension}延拓到 \(L\)。此延拓固定 \(K\)，由正规嵌入判据其像等于 \(L\)，所以属于 \(G_E\)，却不固定 \(a\)。这说明不动域中不可能多出这样的 \(a\)。同一论证也表明，对任何中间域 \(E\)，\(L/E\) 仍正规且可分。

最后取闭子群 \(H\)。显然 \(H\subseteq\operatorname{Gal}(L/L^H)\)。若 \(\sigma\notin H\)，开集 \(G\setminus H\) 含 \(\sigma\)，故有有限 Galois 层 \(N\)，使 \(\sigma U_N\cap H=\varnothing\)。于是 \(\pi_N(\sigma)\notin H_N\coloneqq\pi_N(H)\)。对有限 Galois 扩张 \(N/K\) 应用\cref{b1:thm:finite-galois-correspondence}，有 \(H_N=\operatorname{Gal}(N/N^{H_N})\)。因此存在 \(a\in N^{H_N}\) 不被 \(\pi_N(\sigma)\) 固定。每个 \(h\in H\) 在 \(N\) 上的限制属于 \(H_N\)，所以 \(a\in L^H\)，而 \(\sigma(a)\ne a\)。故 \(\sigma\notin\operatorname{Gal}(L/L^H)\)，得到反向包含。两个操作于是互逆。
\end{proof}
闭性在最后一段有明确用途：它使一个不在 \(H\) 中的自同构能够在某个有限商里与 \(H\) 区分开。没有闭性时，这种有限层分离可能失败。

\subsection{闭性条件与反例}
\label{b1:subsec:2203:02}
子集 \(H\) 的闭包 \(\overline H\) 是包含它的全部闭集之交；一个点属于闭包，当且仅当它的每个开邻域都与 \(H\) 相交。后一个刻画来自取开补集。子群的闭包仍为子群：若 \(x,y\) 的每个有限层邻域都遇到 \(H\)，则可取 \(h,k\in H\) 在同一有限层上分别与 \(x,y\) 相同，\(hk^{-1}\) 就在该层上与 \(xy^{-1}\) 相同，故 \(xy^{-1}\in\overline H\)。

\begin{corollary}\label{b1:cor:fixed-field-closure}
设 \(L/K\) 为代数 Galois 扩张，\(G=\operatorname{Gal}(L/K)\)，并在 \(G\) 上赋予 Krull 拓扑。对任意子群 \(H\leq G\)，记 \(\overline H\) 为 \(H\) 在 \(G\) 中的闭包。则
\[
L^H=L^{\overline H},\qquad
\operatorname{Gal}(L/L^H)=\overline H.
\]
\end{corollary}
\begin{proof}
元素 \(a\) 的稳定子群在上一证明中已经说明是闭子群。因此若它包含 \(H\)，就包含 \(\overline H\)，得到 \(L^H\subseteq L^{\overline H}\)；反向由 \(H\subseteq\overline H\) 得到。对闭子群 \(\overline H\) 应用\cref{b1:thm:infinite-galois-correspondence}，即得第二式。
\end{proof}
下面给出两个不同子群具有相同不动域的例子。取 \(L=\overline{\mathbb F}_q\)、\(K=\mathbb F_q\)，并令 \(F(a)=a^q\)。\cref{b1:thm:finite-field-subfields}及\cref{b1:thm:finite-field-automorphisms}说明有限层为 \(\mathbb F_{q^n}\)，其 Galois 群由 \(F\) 的限制生成，阶为 \(n\)。于是循环子群 \(H=\langle F\rangle\) 在每个有限商中都有满像，故每个基本开陪集都与 \(H\) 相交，即 \(H\) 稠密。

\ProseParagraphBreak
这个循环子群却不等于整个 \(G\)。将 \(n\) 写为 \(n=2^r m\)、\(m\) 为奇数，用中国剩余定理选唯一剩余类 \(e_n\pmod n\)，使
\[
e_n\equiv0\pmod{2^r},\qquad e_n\equiv1\pmod m.
\]
若 \(d\mid n\)，\(e_n\) 模 \(d\) 后满足定义 \(e_d\) 的同一组条件，所以 \(e_n\equiv e_d\pmod d\)。因此有限层上的 \(F^{e_n}\) 相容，\cref{b1:prop:galois-compatible-family}给出一个整体自同构 \(\sigma\)。若它等于某个整数幂 \(F^a\)，则对所有 \(r\) 都有 \(a\equiv0\pmod{2^r}\)，迫使整数 \(a=0\)；但 \(e_3\equiv1\pmod3\)，又排除了 \(a=0\)。故 \(H\ne G\)。

\ProseParagraphBreak
\(F\) 固定的元素满足 \(x^q=x\)，恰为 \(\mathbb F_q\)。所以真子群 \(H\) 与整个 \(G\) 的不动域都是 \(\mathbb F_q\)，说明闭性不能从无限 Galois 对应中删除。

\subsection{有限子扩张、开子群与正规子群}
\label{b1:subsec:2203:03}
\begin{proposition}\label{b1:prop:finite-open-correspondence}
设 \(L/K\) 为代数 Galois 扩张，\(G=\operatorname{Gal}(L/K)\)，并在 \(G\) 上赋予 Krull 拓扑。对中间域 \(K\subseteq E\subseteq L\)，记 \(G_E=\operatorname{Gal}(L/E)\)。则 \([E:K]<\infty\) 当且仅当 \(G_E\) 是 \(G\) 的开子群；此时
\[
[E:K]=[G:G_E].
\]
\end{proposition}
\begin{proof}
若 \(E/K\) 有限，取有限个域生成元并包含在有限 Galois 层 \(N\) 中，则 \(U_N\subseteq G_E\)，所以 \(G_E\) 开。反过来，若 \(G_E\) 开，\cref{b1:thm:krull-compact}给出 \(U_N\subseteq G_E\)。取不动域并用\cref{b1:thm:infinite-galois-correspondence}，得 \(E\subseteq L^{U_N}=N\)，故扩张有限。

选择这样的 \(N\)，限制映射给出陪集集合之间的双射
\(G/G_E\rightarrow G_N/\operatorname{Gal}(N/E)\)：满射来自限制的满射性，两个限制落在同一陪集当且仅当两个原自同构之商固定 \(E\)，所以单射。有限 Galois 对应给出右端的大小为 \([E:K]\)，完成次数公式。
\end{proof}
任意开子群都指数有限。反向若再假设子群闭，则其有限多个陪集都闭，补集也闭，因而子群开。对中间域所对应的子群，\cref{b1:thm:infinite-galois-correspondence}保证闭性，所以有限指数与开性等价。

\begin{theorem}\label{b1:thm:infinite-galois-quotient}
设 \(L/K\) 为代数 Galois 扩张，\(G=\operatorname{Gal}(L/K)\)，并在 \(G\) 上赋予 Krull 拓扑。设 \(H\leq G\) 为闭子群，并令 \(E=L^H\)。则 \(H\) 正规当且仅当 \(E/K\) 正规；在这两个等价条件成立时，\(E/K\) 为 Galois 扩张，限制映射诱导群同构
\[
G/H\cong\operatorname{Gal}(E/K).
\]
\end{theorem}
\begin{proof}
若 \(E/K\) 正规，任何 \(\sigma\in G\) 都保持 \(E\)，故限制给出同态 \(G\rightarrow\operatorname{Gal}(E/K)\)，核为 \(G_E=H\)，所以 \(H\) 正规。任取目标群中的自同构，由\cref{b1:thm:embedding-extension}延拓到 \(L\) 的嵌入，再用 \(L/K\) 正规性得到 \(L\) 的自同构，证明限制满射。群的第一同构定理给出所列同构。

反过来，若 \(H\) 正规，对 \(a\in E,h\in H,\sigma\in G\)，有
\[
h\bigl(\sigma(a)\bigr)
=\sigma\bigl((\sigma^{-1}h\sigma)(a)\bigr)=\sigma(a),
\]
因为 \(\sigma^{-1}h\sigma\in H\)。所以 \(\sigma(E)\subseteq E\)，再用 \(\sigma^{-1}\) 得到相等。任意 \(K\)-嵌入 \(E\rightarrow\overline L\) 可延拓到 \(L\)，延拓属于 \(G\)，因此将 \(E\) 映为自身。\cref{b1:thm:normal-embedding-criterion}说明 \(E/K\) 正规；可分性则从 \(L/K\) 继承。
\end{proof}
这个同构还保持自然拓扑。为了把商拓扑的含义与有限商检测放在一起，我们将在\subsecref{b1:subsec:2205:04}给出其拓扑方面的证明；本小节已完成群与域的代数对应。

% !TeX root = ../../Book1.tex
% 纲要标题：### 22.4 典型绝对 Galois 群
% 纲要位置：计划纲要.md，第 789 行。

\section{典型绝对 Galois 群}
\label{b1:sec:2204}

绝对 Galois 群汇集了基域全部有限可分扩张的信息，但一般很难完整描述。这里有限 Galois 扩张与一般有限可分扩张所对应的群论对象有所不同：有限 Galois 扩张 \(N/K\) 给出有限商群 \(G_K/\operatorname{Gal}(K^{\mathrm{sep}}/N)\cong\operatorname{Gal}(N/K)\)；一般有限可分扩张 \(E/K\) 对应开子群 \(\operatorname{Gal}(K^{\mathrm{sep}}/E)\)，它未必正规，因而未必给出商群。若要用有限商群记录 \(E\)，可先取它的有限正规闭包。上述区别来自\cref{b1:prop:finite-open-correspondence,b1:thm:infinite-galois-quotient}。有限域给出可显式写成逆极限的模型，实数域给出有限模型。它们也说明，连续到有限群的映射必须能由某个有限层决定。

% 纲要内容（仅作写作计划，尚非教材正文）：
%
% * 有限域的绝对 Galois 群；区分全部模数的投射有限整数与单一素数的进整数，以及拓扑生成与抽象生成
% * 实数域与实闭域的联系
% * 连续表示与离散集上的连续作用；正式证明到有限离散群的连续同态由整数 1 的像唯一决定，再解释任意可逆矩阵均可指定为 Frobenius 的像；说明连续群上同调的后续位置
%

\subsection{有限域的绝对 Galois 群}
\label{b1:subsec:2204:01}
把正整数按整除关系排列；当 \(n\mid m\) 时，以模 \(n\) 的约化
\(\mathbb Z/m\mathbb Z\rightarrow\mathbb Z/n\mathbb Z\) 为过渡映射。所得拓扑环
\[
\widehat{\vphantom{Z}\mathbb Z}=\varprojlim_{n\geq1}\mathbb Z/n\mathbb Z,
\]
称为\term[toushe-youxian-zhengshu]{投射有限整数环}[profinite integers]。其元素是相容剩余类族 \((a_n)_n\)，加法与乘法均逐坐标进行；拓扑是有限离散环直积中的子空间拓扑。
\symidx{\(\widehat{\vphantom{Z}\mathbb Z}\)：投射有限整数环}

\begin{theorem}\label{b1:thm:finite-field-absolute-galois}
设 \(q\) 为素数幂，\(G_{\mathbb F_q}=\operatorname{Gal}(\overline{\mathbb F}_q/\mathbb F_q)\) 为有限域 \(\mathbb F_q\) 的绝对 Galois 群，并在其上赋予 Krull 拓扑。令 \(\widehat{\vphantom{Z}\mathbb Z}=\varprojlim_{n\geq1}\mathbb Z/n\mathbb Z\) 按整除关系及自然约化构成，并取逆极限拓扑。则存在拓扑群同构
\[
\widehat{\vphantom{Z}\mathbb Z}\longrightarrow G_{\mathbb F_q},
\qquad(a_n)_n\longmapsto\sigma,
\qquad \sigma|_{\mathbb F_{q^n}}=F^{a_n},
\]
其中 \(F(x)=x^q\)。整数 \(1\) 对应 Frobenius 自同构。
\end{theorem}
\begin{proof}
有限域完美，其可分闭包就是代数闭包。每个代数元 \(a\in\overline{\mathbb F}_q\) 都生成有限扩张 \(\mathbb F_q(a)/\mathbb F_q\)；若其次数为 \(n\)，这个域便有 \(q^n\) 个元素，因此是当前闭包中由 \(X^{q^n}-X\) 的全部根组成的唯一子域 \(\mathbb F_{q^n}\)。由\cref{b1:thm:finite-field-subfields}，这些有限子域的包含关系对应 \(n\mid m\)。故它们已经覆盖整个代数闭包，并给出全部有限 Galois 层。

由\cref{b1:thm:finite-field-automorphisms}，第 \(n\) 层的自同构群是阶 \(n\) 的循环群，以 \(F\) 的限制为生成元。把它与加法群 \(\mathbb Z/n\mathbb Z\) 识别以后，\(F^a\) 从第 \(m\) 层限制到第 \(n\) 层，只需把指数 \(a\) 模 \(n\) 约化。因此\cref{b1:prop:galois-compatible-family}给出所述群同构。在逆极限一侧，规定第 \(n\) 个坐标的剩余类，对应于在 Galois 群一侧规定向 \(\mathbb F_{q^n}\) 的限制；有限多个坐标条件又可合并到一个共同倍数所对应的坐标。这两类条件分别构成两侧拓扑的邻域基，故该群同构及其逆都连续。
\end{proof}
整数到 \(\widehat{\vphantom{Z}\mathbb Z}\) 的自然映射为 \(a\mapsto(a\bmod n)_n\)，它单射，因为被所有正整数整除的整数只能为零。其像稠密：有限多个模数的条件可由一个共同倍数的坐标统一规定，再选该剩余类的一个整数代表。\subsecref{b1:subsec:2203:02}中构造的 \((e_n)_n\) 表明这个像不等于全部投射有限整数。因此整数幂组成的普通循环子群 \(\langle F\rangle=\{F^a:a\in\mathbb Z\}\cong\mathbb Z\) 是真子群，却满足
\[
\overline{\langle F\rangle}=G_{\mathbb F_q}.
\]
称 \(F\) 为拓扑生成元，是说它的整数幂稠密，而不是说每个群元素都是某个整数幂。

\ProseParagraphBreak
\(\widehat{\vphantom{Z}\mathbb Z}\) 同时记录所有正整数模数的相容信息。给定素数 \(p\)，改为只记录模数为 \(p^r\) 的坐标，所得环
\(\mathbb Z_p=\varprojlim_r\mathbb Z/p^r\mathbb Z\) 称为 \(p\) 进整数环；从 \(\widehat{\vphantom{Z}\mathbb Z}\) 到 \(\mathbb Z_p\) 的自然映射就是舍去其余坐标。两种逆极限所保留的有限信息不同。以后讨论 \(p\) 进表示时，正是这些模 \(p^r\) 的有限信息决定连续性。
\symidx{\(\mathbb Z_p\)：\(p\) 进整数环}

\subsection{实数域与实闭域}
\label{b1:subsec:2204:02}
\begin{proposition}\label{b1:prop:real-absolute-galois}
实数域的绝对 Galois 群为 \(G_{\mathbb R}=\operatorname{Gal}(\mathbb C/\mathbb R)\cong C_2\)，其中非单位元为复共轭；其 Krull 拓扑为离散拓扑。
\end{proposition}
\begin{proof}
由 \(\mathbb C=\mathbb R(i)\)、\(i^2=-1\)，以及 \(X^2+1\) 在 \(\mathbb R\) 上不可约，得 \([\mathbb C:\mathbb R]=2\)，所以 \(\mathbb C/\mathbb R\) 是代数扩张。\cref{b1:thm:fundamental-theorem-algebra}又保证 \(\mathbb C\) 代数闭，故它确为 \(\mathbb R\) 的一个代数闭包；该定理的证明见\secref{b1:sec:E02}。特征零保证所有代数元可分，所以这个代数闭包也是可分闭包。

\(\mathbb C\) 是 \(X^2+1\) 的分裂域。自同构必须把 \(i\) 送到 \(i\) 或 \(-i\)，两种选择分别给出恒等与复共轭，所以群为 \(C_2\)。整个闭包已经是有限 Galois 层，因而 \(\operatorname{Gal}(\mathbb C/\mathbb C)=\{1\}\) 是基本开子群；它的每个陪集都是单点开集，故拓扑离散。
\end{proof}
这个例子的代数闭性来自实闭域刻画与实分析的中间值定理；自同构群及其拓扑则由二次扩张直接算出。

\ProseParagraphBreak
相同计算对特征零、满足“\(-1\) 不是平方且添加 \(i\) 后得到代数闭域”的域 \(F\) 也适用：\(F(i)/F\) 为二次可分闭包，故 \(G_F\cong C_2\)。这些条件是实闭域的一种代数刻画，具体见\cref{b1:thm:real-closed-characterizations}；\appref{b1:app:E}从有序域与实闭包出发加以说明。此处只作这一条件下的计算，不把任意有序域或任意实数子域都视为实闭域。例如 \(\mathbb Q\subseteq\mathbb R\)，但 \(\mathbb Q\) 还有许多不同次数的有限 Galois 扩张。

\subsection{连续表示与离散集上的连续作用}
\label{b1:subsec:2204:03}
\begin{definition}\label{b1:def:continuous-finite-representation}
设 \(K\) 为域，在一个固定代数闭包中取可分闭包 \(K^{\mathrm{sep}}\)，并令 \(G_K=\operatorname{Gal}(K^{\mathrm{sep}}/K)\) 为赋予 Krull 拓扑的绝对 Galois 群；再设 \(k\) 为有限域，\(d\geq1\) 为整数。如果群同态
\[
\rho:G_K\longrightarrow\operatorname{GL}_d(k),
\]
在目标群 \(\operatorname{GL}_d(k)\) 取离散拓扑时连续，则称 \(\rho\) 为 \(G_K\) 在 \(k^d\) 上的 \(d\) 维\term[lianxu-biaoshi]{连续表示}[continuous representation]。这里 \(\operatorname{GL}_d(k)\) 表示可逆矩阵群，\(\rho\) 规定群元素在线性空间 \(k^d\) 上怎样作用。
\end{definition}
\symidx{\(\operatorname{GL}_d(k)\)：域上的一般线性群}
连续性使 \(\rho^{-1}(\{I_d\})\) 成为单位元的开邻域，所以它包含某个 \(U_N=\operatorname{Gal}(K^{\mathrm{sep}}/N)\)，其中 \(N/K\) 是有限 Galois 扩张。这意味着两个自同构只要在 \(N\) 上限制相同，它们作用在 \(k^d\) 上的矩阵就相同：每个连续表示的全部作用信息已经由某个有限 Galois 商决定。这个判别及其逆命题将在\cref{b1:prop:finite-continuous-hom}中完整证明。对于有限域，Frobenius 的一个像就能确定这样的表示，但这里还需要连续性。

\begin{proposition}\label{b1:prop:profinite-integer-finite-hom}
设 \(D\) 为赋予离散拓扑的有限群，并令 \(\widehat{\vphantom{Z}\mathbb Z}=\varprojlim_{n\geq1}\mathbb Z/n\mathbb Z\) 按整除关系及自然约化构成，赋予逆极限拓扑。把 \(\widehat{\vphantom{Z}\mathbb Z}\) 看成加法群，并用 \(1\) 表示整数 \(1\) 的自然像。则连续群同态
\[
\rho:\widehat{\vphantom{Z}\mathbb Z}\longrightarrow D
\]
通过 \(\rho\mapsto\rho(1)\) 与 \(D\) 的元素建立一一对应。若 \(d\in D\) 的阶为 \(m\)，对应同态可先约化到 \(\mathbb Z/m\mathbb Z\)，再令 \(\bar a\mapsto d^a\)。
\end{proposition}
\begin{proof}
给定 \(d\in D\)，它的阶 \(m\) 有限。因为 \(d^m=1_D\)，映射 \(\mathbb Z/m\mathbb Z\rightarrow D\)、\(\bar a\mapsto d^a\) 定义良好且为群同态。把它与第 \(m\) 个坐标投影复合，便得到所需同态。坐标投影连续，有限离散群之间的映射也连续，所以这个复合同态连续。

若两个连续同态 \(\rho,\eta\) 在 \(1\) 上取值相同，则由同态性，它们在所有整数的自然像上取值相同。这个整数子群稠密，而等值集合
\[
\{x:\rho(x)=\eta(x)\}
=\bigcup_{c\in D}\bigl(\rho^{-1}(\{c\})\cap\eta^{-1}(\{c\})\bigr)
\]
是闭集：目标离散，单点闭，两映射连续，并且这里是有限个闭集的并。等值集合包含稠密子群，故为整个 \(\widehat{\vphantom{Z}\mathbb Z}\)，这就证明唯一性。
\end{proof}
这里的唯一性来自整数子群的稠密性与同态的连续性，不能用“\(1\) 生成整个抽象群”来解释。由\cref{b1:thm:finite-field-absolute-galois}，对任一 \(A\in\operatorname{GL}_d(k)\)，恰有一个连续表示 \(\rho:G_{\mathbb F_q}\rightarrow\operatorname{GL}_d(k)\) 满足 \(\rho(F)=A\)；它经过模 \(\operatorname{ord}(A)\) 的有限循环商。这是有限域绝对 Galois 群的特殊性质，一般绝对 Galois 群的元素及其关系会限制可指定的矩阵像。

\ProseParagraphBreak
例如设 \(\ell\) 为素数，取系数域 \(k=\mathbb F_\ell\)。矩阵
\[
A=\begin{pmatrix}1&1\\0&1\end{pmatrix}
\]
满足 \(A^r=\begin{psmallmatrix}1&r\\0&1\end{psmallmatrix}\)，因此其阶为 \(\ell\)。令 Frobenius 元 \(F\) 的像为 \(A\)，便得到经过 \(\mathbb Z/\ell\mathbb Z\) 的非平凡二维连续表示。每个群元素的像都是某个 \(A^r\)，其迹均为 \(2\cdot1_k\)，与二维平凡表示的迹相等；两表示却不同构，因为 \(F\) 在其中分别作用为 \(A\) 与单位矩阵，而 \(A\ne I_2\) 不可能与 \(I_2\) 共轭。

\ProseParagraphBreak
这个表示中，由第一个标准基向量张成的直线是平凡子表示，商空间上的作用也平凡，但整个表示不是两个平凡表示的直和。迹只给出这两个一维作用的迹之和，不能辨认它们在这里怎样组合成二维作用。系数域的特征 \(\ell\) 正好整除有限像群的阶 \(\ell\)，这一例说明在这种特征下，迹函数可能无法区分不同构的表示。

\ProseParagraphBreak
设拓扑群 \(G_K\) 作用在离散集合 \(M\) 上。如果作用映射 \(G_K\times M\rightarrow M\) 连续，则称这个作用连续。在 \(M\) 离散时，这等价于每个元素的稳定子群开：连续性在 \((1,m)\) 附近给出一个固定 \(m\) 的开邻域；反向则在每个 \((g,m)\) 用稳定子群的陪集构造保持作用值的开邻域。第四卷附录C将定义连续离散模的群上同调，并介绍相应的上链与系数条件；本章的拓扑群知识是学习这些内容的基础。

% !TeX root = ../../Book1.tex
% 纲要标题：### 22.5 有限层检测、闭包与连续性
% 纲要位置：计划纲要.md，第 760 行。

\section{有限层检测、闭包与连续性}
\label{b1:sec:2205}

闭子群可以通过全部有限商中的像辨认。到有限离散群的同态，其连续性可以通过核是否包含某个基本开子群判定；到赋有逆极限拓扑的群的同态，则可逐个有限坐标检验连续性。本节将这些判别写成公式，并说明商群的拓扑为何仍须保留闭性要求。

% 纲要内容（仅作写作计划，尚非教材正文）：
%
% * 闭子群由其在各有限 Galois 商中的像确定；写出子群闭包的有限商检测公式
% * 以有限域的 Frobenius 生成的稠密循环子群比较 \(\mathbb Z\) 与 \(\widehat{\mathbb Z}\)，说明具有同一不动域的子群为何必须取闭包
% * 连续到有限离散群的同态具有开核；有限层上的表示与相容的连续表示
% * 无限 Galois 对应中的商群结论要求闭正规子群；说明与第四卷连续上同调的联系，不调用其定理
%
% ---
%

\subsection{闭子群与闭包的有限商检测}
\label{b1:subsec:2205:01}
\begin{proposition}\label{b1:prop:closure-finite-quotients}
设 \(L/K\) 为代数 Galois 扩张，\(G=\operatorname{Gal}(L/K)\)，并在 \(G\) 上赋予 Krull 拓扑。令 \(\mathcal N\) 为 \(L/K\) 的全部有限 Galois 子扩张组成的有向集合；对 \(N\in\mathcal N\)，记 \(\pi_N:G\rightarrow\operatorname{Gal}(N/K)\) 为限制映射，\(U_N=\ker\pi_N=\operatorname{Gal}(L/N)\)。则对任意子群 \(H\leq G\)，有
\begin{equation}\label{b1:eq:closure-finite-quotients}
\overline H=\bigcap_{N\in\mathcal N}\pi_N^{-1}\bigl(\pi_N(H)\bigr)
=\bigcap_{N\in\mathcal N}HU_N.
\end{equation}
其中 \(\overline H\) 表示 \(H\) 在 \(G\) 中的闭包。所以两个闭子群相等，当且仅当它们在每个有限 Galois 商中的像相等。
\end{proposition}
\begin{proof}
点 \(\sigma\) 属于 \(\overline H\)，当且仅当每个基本开邻域 \(\sigma U_N\) 与 \(H\) 相交。相交意味着有 \(h\in H\) 满足 \(\pi_N(h)=\pi_N(\sigma)\)，这等价于 \(\sigma\in\pi_N^{-1}\bigl(\pi_N(H)\bigr)\)，也等价于 \(\sigma\in HU_N\)。逐个 \(N\) 求交得到公式。若两个闭子群的有限像相同，两者分别等于公式中的同一个交；反向由取像直接成立。
\end{proof}
把这个公式与\cref{b1:cor:fixed-field-closure}合用，得到一个完整的辨别准则：任意两个子群具有相同不动域，当且仅当它们的闭包相同，也当且仅当它们在全部有限商中的像相同。对子群本身而言，有限信息只确定它的闭包。

\subsection{Frobenius 的幂、闭包与有限子域}
\label{b1:subsec:2205:02}
设 \(q\) 为素数幂，\(F(x)=x^q\)。在 \(G_{\mathbb F_q}=\operatorname{Gal}(\overline{\mathbb F}_q/\mathbb F_q)\cong\widehat{\vphantom{Z}\mathbb Z}\) 中，\(\langle F\rangle\) 对应整数的自然像 \(\mathbb Z\)。\subsecref{b1:subsec:2203:02}已证明这一像稠密而非全群；现在考虑 \(F\) 的幂所生成的子群。

\ProseParagraphBreak
固定正整数 \(d\)，有
\begin{equation}\label{b1:eq:profinite-multiple-kernel}
d\widehat{\vphantom{Z}\mathbb Z}
=\ker\left(\widehat{\vphantom{Z}\mathbb Z}\longrightarrow\mathbb Z/d\mathbb Z\right).
\end{equation}
一个方向由乘 \(d\) 直接得到。反过来，设相容族 \(a\) 在模 \(d\) 下为零。对每个 \(n\)，选择 \(a_{dn}\) 的一个被 \(d\) 整除的整数代表，除以 \(d\) 后取模 \(n\)，记为 \(b_n\)。改变代表只改变 \(n\) 的倍数，故定义不变；原族的相容性说明 \((b_n)_n\) 相容，并满足 \(db=a\)。这证明另一方向。

\ProseParagraphBreak
\(\langle F^d\rangle\) 在第 \(n\) 个商 \(\mathbb Z/n\mathbb Z\) 中的像为
\[
d(\mathbb Z/n\mathbb Z)=\{[dk]_n:k\in\mathbb Z\}
=\frac{d\mathbb Z+n\mathbb Z}{n\mathbb Z}.
\]
这些有限像也正是 \(d\widehat{\vphantom{Z}\mathbb Z}\) 的有限像；后者由\eqref{b1:eq:profinite-multiple-kernel}知是闭子群。因此\eqref{b1:eq:closure-finite-quotients}给出
\(\overline{\langle F^d\rangle}=d\widehat{\vphantom{Z}\mathbb Z}\)。它恰为固定 \(\mathbb F_{q^d}\) 的自同构组成的子群，所以 \(\langle F^d\rangle\) 与其闭包的不动域同为 \(\mathbb F_{q^d}\)。

下表比较这几个子群的拓扑性质与不动域，其中 \(d\geq2\)，\(\mathbb Z\) 表示整数在 \(\widehat{\vphantom{Z}\mathbb Z}\) 中的自然像。\(d\widehat{\vphantom{Z}\mathbb Z}\) 由\eqref{b1:eq:profinite-multiple-kernel}知既闭又开；\(\mathbb Z\) 稠密而非全群，因而不闭，也不可能开。零子群闭；若它开，紧群 \(\widehat{\vphantom{Z}\mathbb Z}\) 就成为有限离散空间，与它无限相矛盾。特别地，\(\mathbb Z\) 与全群虽不相等，却有同一个不动域。
\begin{center}
\begin{tabular}{@{}lccc@{}}
\multicolumn{4}{c}{Frobenius 子群的闭性、开性与不动域}\\
\addlinespace[0.4em]
\toprule
子群 \(H\) & 闭 & 开 & \(\overline{\mathbb F}_q^{\,H}\) \\
\midrule
\(\widehat{\vphantom{Z}\mathbb Z}\) & 是 & 是 & \(\mathbb F_q\) \\
\(d\widehat{\vphantom{Z}\mathbb Z}\) & 是 & 是 & \(\mathbb F_{q^d}\) \\
\(\mathbb Z\) & 否 & 否 & \(\mathbb F_q\) \\
\(\{0\}\) & 是 & 否 & \(\overline{\mathbb F}_q\) \\
\bottomrule
\end{tabular}
\end{center}

\subsection{开核与有限层表示}
\label{b1:subsec:2205:03}
\begin{proposition}\label{b1:prop:finite-continuous-hom}
设 \(L/K\) 为代数 Galois 扩张，\(G=\operatorname{Gal}(L/K)\)，并在 \(G\) 上赋予 Krull 拓扑。对每个有限 Galois 子扩张 \(N/K\)，记 \(G_N=\operatorname{Gal}(N/K)\)，\(\pi_N:G\rightarrow G_N\) 为限制映射。再设 \(D\) 为有限离散群，\(\rho:G\rightarrow D\) 为群同态。以下条件等价：\(\rho\) 连续；\(\ker\rho\) 开；存在某个这样的 \(N\) 及同态 \(\rho_N:G_N\rightarrow D\)，使 \(\rho=\rho_N\pi_N\)。
\end{proposition}
\begin{proof}
连续性使单点 \(\{1\}\) 的原像开。若核开，取 \(U_N\subseteq\ker\rho\)，规定 \(\rho_N(\sigma|_N)=\rho(\sigma)\)。两个提升之商在 \(U_N\) 中，故像相同；限制满射保证公式在全部 \(G_N\) 上有定义，并保持群乘法。反过来，有限层投影连续，有限离散集合之间的任意映射连续，所以经过有限商的同态连续。
\end{proof}
有限层不一定唯一。若某表示经过 \(G_N\)，也经过任何包含 \(N\) 的更大有限 Galois 层；应比较其限制相容性，而不是把选择层的不同误当作表示不同。

\ProseParagraphBreak
取素数 \(\ell\) 和整数 \(d\geq1\)。相容的有限层表示可合成以 \(\ell\) 进整数环为系数的表示。给 \(M_d(\mathbb Z_\ell)\cong\mathbb Z_\ell^{d^2}\) 赋予积拓扑，并给 \(\operatorname{GL}_d(\mathbb Z_\ell)\) 赋予子空间拓扑；这里的矩阵群不取离散拓扑。作为拓扑群，\(\operatorname{GL}_d(\mathbb Z_\ell)\) 可识别为
\[
\varprojlim_r\operatorname{GL}_d(\mathbb Z/\ell^r\mathbb Z).
\]
相容矩阵逐项给出 \(\mathbb Z_\ell\) 中的矩阵，有限层的逆矩阵也相容，从而给出整体逆矩阵。两侧拓扑均由模 \(\ell^r\) 的坐标条件生成，故这一识别还是同胚。具体地，若
\[
V_r=\ker\!\left(\operatorname{GL}_d(\mathbb Z_\ell)
\longrightarrow\operatorname{GL}_d(\mathbb Z/\ell^r\mathbb Z)\right),
\]
则 \(AV_r\)（\(r\geq1\)）构成矩阵 \(A\in\operatorname{GL}_d(\mathbb Z_\ell)\) 的邻域基。若 \(L/K\) 为代数 Galois 扩张，\(G=\operatorname{Gal}(L/K)\) 取 Krull 拓扑，且 \(\rho_r:G\rightarrow\operatorname{GL}_d(\mathbb Z/\ell^r\mathbb Z)\) 均连续、模约化后相容，就得到唯一同态 \(\rho:G\rightarrow\operatorname{GL}_d(\mathbb Z_\ell)\)。目标的基本开集只限制有限个坐标，其原像为有限多个开集的交，故 \(\rho\) 连续；反向将连续的 \(\rho\) 与各坐标投影复合即可。每个 \(\rho_r\) 都可由\cref{b1:prop:finite-continuous-hom}经过某个有限 Galois 商 \(\operatorname{Gal}(N_r/K)\)。这里 \(r\) 记录矩阵系数的精度，\(N_r\) 记录决定这一级表示的域扩张，两者是不同的层次；随着 \(r\) 增大，所需的 \(N_r\) 也可以增大。逐个精度能经过有限商，并不要求存在一个固定有限层同时决定全部精度，因而整体 \(\rho\) 不必有开核或有限像。

\ProseParagraphBreak
以有限域为例，由\cref{b1:thm:finite-field-absolute-galois}，可以把 \(G_{\mathbb F_q}\) 识别为加法群 \(\widehat{\vphantom{Z}\mathbb Z}\)。对相容族 \(a=(a_n)_n\)，只取素数幂模数的坐标，得到
\[
a_\ell=(a_{\ell^r})_{r\geq1}\in\mathbb Z_\ell.
\]
定义二维表示
\[
\rho(a)=\begin{pmatrix}1&a_\ell\\0&1\end{pmatrix}
\in\operatorname{GL}_2(\mathbb Z_\ell).
\]
矩阵乘法将右上角相加，故这是群同态。其模 \(\ell^r\) 约化只依赖 \(a_{\ell^r}\)，因此连续，并经过有限 Galois 层 \(N_r=\mathbb F_{q^{\ell^r}}\)；全部约化相容，所以 \(\rho\) 连续。普通整数 \(k\) 的像是 \(\begin{psmallmatrix}1&k\\0&1\end{psmallmatrix}\)，不同整数在 \(\mathbb Z_\ell\) 中仍不同，故整体像无限。若所有约化都经过同一个有限商，那么两元素在该商中相等时，其矩阵在每个模 \(\ell^r\) 下都相等，从而整体矩阵相等；这会使 \(\rho\) 本身经过该有限商，像也有限，矛盾。因此这一例中没有统一的有限层。其核也不开：若核开，紧性使其陪集只有有限个，而这些陪集与表示的像一一对应，又会迫使像有限。

\subsection{闭正规子群与 Galois 商群}
\label{b1:subsec:2205:04}
设 \(G\) 为拓扑群，\(H\trianglelefteq G\) 为正规子群，\(q:G\rightarrow G/H\) 为自然映射。在抽象商群 \(G/H\) 上定义\term[shang-tuopu]{商拓扑}[quotient topology]：集合 \(V\subseteq G/H\) 开，当且仅当 \(q^{-1}(V)\) 在 \(G\) 中开。本节应用时，\(G=\operatorname{Gal}(L/K)\)，\(H\) 还是闭子群，并令 \(E=L^H\)。

\begin{proposition}\label{b1:prop:galois-quotient-homeomorphism}
设 \(L/K\) 为代数 Galois 扩张，\(G=\operatorname{Gal}(L/K)\)，并在 \(G\) 上赋予 Krull 拓扑。设 \(H\leq G\) 为闭正规子群，\(E=L^H\)。在 \(G/H\) 上赋予商拓扑，在 \(\operatorname{Gal}(E/K)\) 上赋予 Krull 拓扑。则限制映射诱导的群同构
\(G/H\rightarrow\operatorname{Gal}(E/K)\) 是拓扑群同构。
\end{proposition}
\begin{proof}
先看限制同态 \(r:G\rightarrow\operatorname{Gal}(E/K)\)。目标中一个基本邻域规定某个有限 Galois 子扩张 \(F/K\)、\(F\subseteq E\) 上的作用，其原像在 \(G\) 中也是同样规定 \(F\) 上作用的基本开集，所以 \(r\) 连续。写成 \(r=\bar r q\)，由商拓扑定义可知诱导双射 \(\bar r\) 连续。

若 \(C\subseteq G/H\) 闭，则 \(q^{-1}(C)\) 是紧群 \(G\) 的闭子集，因而紧。由\cref{b1:prop:compact-maps}，其连续像 \(r\bigl(q^{-1}(C)\bigr)=\bar r(C)\) 在 Hausdorff 群 \(\operatorname{Gal}(E/K)\) 中紧，故闭。因此 \(\bar r\) 还是闭映射，逆映射连续，所列同构为同胚。
\end{proof}
若 \(H\) 是非闭正规子群，\(G/H\) 的商拓扑就不可能是 Hausdorff：否则单位元单点 \(\{H\}\) 闭，由商映射 \(q\) 的连续性，\(H=q^{-1}(\{H\})\) 也闭，矛盾。以 \(G=\widehat{\vphantom{Z}\mathbb Z}\)、\(H=\mathbb Z\) 的稠密像为例，商群非平凡，商拓扑却只有空集和全空间两个开集。事实上，若商中有非空开集，其原像 \(V\) 是非空开集且满足 \(V+H=V\)。对任意 \(g\in G\)，开集 \(g-V\) 与稠密的 \(H\) 相交，于是 \(g\in V+H=V\)，故 \(V=G\)。

\ProseParagraphBreak
无限 Galois 对应将中间域与闭子群配对，有限中间扩张对应开子群，正规中间扩张对应闭正规子群；相应中间扩张的 Galois 群是原群对该闭正规子群的拓扑商群。更进一步的连续群上同调和算术表示，还需指定系数对象及相应的连续性条件。


\begin{chaptersummary}
有限层上的限制将一个 Galois 自同构变为相容坐标族，Krull 拓扑则规定有限个坐标所能观察到的信息。相容条件把逆极限刻画为有限离散群直积中的闭子群，由此得到紧性、Hausdorff 性和全不连通性。证明中的选择原理用于极大真滤子的存在。

\ProseParagraphBreak
闭子群恰好能从其不动域恢复。任意子群与其闭包有相同不动域，有限域中Frobenius生成的稠密真子群给出这一现象的明确例子。有限子扩张对应开子群，正规子扩张对应闭正规子群；商群同构还须带上自然的商拓扑。

\ProseParagraphBreak
这些结论使无限对象仍可通过有限层计算。连续到有限离散群的同态经过某个有限Galois商，相容的有限层表示则可以合成到逆极限系数中的连续表示。有限与无限Galois理论由此有了统一的描述；无限情形须要求子群闭、同态连续。进一步的算术与上同调应用将以这些结论为基础。
\end{chaptersummary}

\BookExercises

\begin{exercise}\label{b1:ex:22-1}
固定 \(\overline{\mathbb F}_q\)，令 \(E_r=\mathbb F_{q^{r!}}\)。证明 \(\overline{\mathbb F}_q=\bigcup_{r\geq1}E_r\)。若给定 \(\tau_r\in\operatorname{Gal}(E_r/\mathbb F_q)\)，且 \(\tau_{r+1}|_{E_r}=\tau_r\)，证明它们唯一确定 \(\overline{\mathbb F}_q\) 的一个 \(\mathbb F_q\)-自同构，并说明只用阶乘层的限制条件仍给出相同的 Krull 拓扑。
\end{exercise}
\begin{solution}
每个代数元属于某个 \(\mathbb F_{q^n}\)。若 \(r\geq n\)，则 \(n\mid r!\)，由\cref{b1:thm:finite-field-subfields}，\(\mathbb F_{q^n}\subseteq E_r\)，故这些域的并是整个闭包。这里并不要求每个有限域都等于某个 \(E_r\)；关键是每个有限层都包含在某个阶乘层中。给定题中的相容族，对任意有限层 \(\mathbb F_{q^n}\) 选取 \(r\) 使 \(n\mid r!\)，用 \(\tau_r\) 的限制定义其作用。换另一 \(s\) 时再取更大的阶乘层，相容性保证两次限制相同。这样得到全部有限层的相容族，且每个限制都固定 \(\mathbb F_q\)，由\cref{b1:prop:galois-compatible-family}唯一拼成整体 \(\mathbb F_q\)-自同构。阶乘层给出的开陪集本来就是 Krull 基本开集，而任意有限层给出的基本开集又能由更大的阶乘层细化，所以只保留这些层也给出相同的 Krull 拓扑。
\end{solution}

\begin{exercise}\label{b1:ex:22-2}
设 \(K=\mathbb F_p(t)\)，在固定代数闭包中令 \(L=\bigcup_{r\geq0}K(t^{1/p^r})\)。证明 \(L/K\) 无限，且 \(\operatorname{Aut}_K(L)=1\)。为什么不能在此使用无限 Galois 对应？
\end{exercise}
\begin{solution}
在 \(\mathbb F_p[t][X]\) 中，多项式 \(X^{p^r}-t\) 对素元 \(t\) 满足\cref{b1:thm:eisenstein}：常数项被 \(t\) 整除但不被 \(t^2\) 整除，其他非首项系数为零。因此 \(K(t^{1/p^r})/K\) 的次数为 \(p^r\)，这些次数无界，故 \(L/K\) 无限。每个 \(t^{1/p^r}\) 是 \(X^{p^r}-t\) 在闭包中的唯一根，任何 \(K\) 自同构都必须固定它，所以固定整个并域。其自同构群的不动域为 \(L\)，并非 \(K\)。扩张是纯不可分的，不满足\cref{b1:thm:infinite-galois-correspondence}的可分性假设。
\end{solution}

\begin{exercise}\label{b1:ex:22-3}
在 \(\overline{\mathbb F}_q\) 中令 \(F(x)=x^q\)。能否同时在 \(\mathbb F_{q^4}\) 上指定 \(F\)，在 \(\mathbb F_{q^6}\) 上指定 \(F^2\)，再延拓为一个整体自同构？若把后一指定改成 \(F^3\)，答案如何？
\end{exercise}
\begin{solution}
两域的交为 \(\mathbb F_{q^2}\)。第一种指定在交域上的限制分别为非恒等的 \(F\) 与恒等的 \(F^2\)，违反\eqref{b1:eq:galois-compatibility}，所以不可能同时延拓。第二种指定在交域上相同。相应的指数须满足
\[
a\equiv1\pmod4,\qquad a\equiv3\pmod6.
\]
这里模数不互素；可解性要求两个指定指数模 \(\gcd(4,6)=2\) 相同，恰好就是交域上的相容性。写 \(a=1+4u\)，第二个同余式化为 \(2u\equiv1\pmod3\)，即 \(u\equiv2\pmod3\)，所以 \(a\equiv9\pmod{12}\)。整体自同构 \(F^9\) 因而同时实现这两个指定。分别可以延拓不等于所选延拓彼此相容，交域上的限制是必须检查的条件。
\end{solution}

\begin{exercise}\label{b1:ex:22-4}
对每个 \(r\geq1\) 令 \(A_r=\mathbb Z/2\mathbb Z\)，过渡映射 \(A_{r+1}\rightarrow A_r\) 全为零映射。求 \(\varprojlim_r A_r\)，并说明逆极限的坐标投影是否必为满射。
\end{exercise}
\begin{solution}
相容族 \((a_r)_r\) 必须满足 \(a_r=0(a_{r+1})=0\)，故只有零族，逆极限为平凡群。每个坐标投影的像都只有零，而 \(A_r\) 有两个元素，所以不满射。\cref{b1:lem:galois-restriction-surjective}为 Galois 系统提供了额外的满射性，不能把它误认为任意逆极限定义的直接后果。
\end{solution}

\begin{exercise}\label{b1:ex:22-5}
设 \((A_i,r_{ji})\) 为有向指标上的有限群逆系统，且所有过渡映射均满射。证明逆极限到每个 \(A_i\) 的坐标投影满射，明确指出紧性在哪一步使用。
\end{exercise}
\begin{solution}
固定 \(i\) 与 \(a\in A_i\)。在非空积空间 \(X=\prod_j A_j\) 中，要求第 \(i\) 坐标为 \(a\)，并要求所有相容等式成立。这些要求分别定义闭集。任取有限多个要求，选取同时大于固定指标 \(i\) 及这一有限组要求中全部指标的 \(k\)。由 \(r_{ki}\) 满射，取 \(b\in A_k\) 使其在第 \(i\) 层的像为 \(a\)；把所有涉及坐标指定为 \(b\) 的相应像，未涉及坐标取单位元，就满足这一有限组要求。因此闭集族有有限交性质。

\cref{b1:lem:finite-discrete-product-compact}使总交非空，其中任一点都是第 \(i\) 坐标为 \(a\) 的相容族。此处从全部有限组可解推出同时可解，正是紧性步骤；正文给出的紧性证明在极大真滤子的构造处使用 Zorn 引理。
\end{solution}

\begin{exercise}\label{b1:ex:22-6}
令 \(\mathbb Z_2=\varprojlim_r\mathbb Z/2^r\mathbb Z\)。证明乘以 \(2\) 的映射 \(\mathbb Z_2\rightarrow\mathbb Z_2\) 单射，其像恰为模 \(2\) 投影的核，因而是开且闭的指数 \(2\) 子群。
\end{exercise}
\begin{solution}
若 \(2a=0\)，则在模 \(2^{r+1}\) 的坐标上，\(2a_{r+1}=0\)，故 \(a_{r+1}\) 被 \(2^r\) 整除。约化到第 \(r\) 层便得 \(a_r=0\)。对所有 \(r\) 都成立，所以 \(a=0\)。像当然落在模 \(2\) 的核中。

反过来，若 \(a_1=0\)，则每个 \(a_{r+1}\) 有偶数代表。将这个代表除以 \(2\) 再模 \(2^r\)，记为 \(b_r\)；换代表改变的是 \(2^r\) 的倍数，所以良定义。相容性保证 \((b_r)_r\in\mathbb Z_2\)，并有 \(2b=a\)。于是像就是模 \(2\) 的核。该投影满射到离散的二元群，其核开且闭，指数为 \(2\)。
\end{solution}

\begin{exercise}\label{b1:ex:22-7}
设 \(L/K\) 为代数 Galois 扩张，\(G=\operatorname{Gal}(L/K)\) 取 Krull 拓扑。对有限 Galois 中间域 \(N/K\)，记 \(G_N=\operatorname{Gal}(N/K)\)，并记限制映射为 \(\pi_N:G\to G_N\)。证明每个既开又闭的子集 \(C\subseteq G\) 都可写成 \(C=\pi_N^{-1}(D)\)，其中 \(N/K\) 为某个有限 Galois 中间域，\(D\subseteq G_N\)。这里不要求 \(C\) 是子群。
\end{exercise}
\begin{solution}
\(C\) 闭，所以由\cref{b1:thm:krull-compact}及闭子空间的紧性，\(C\) 紧。\(C\) 开，故其每个点有一个包含在 \(C\) 中的基本开陪集。这些陪集覆盖 \(C\)，紧性给出有限子覆盖 \(\sigma_jU_{N_j}\)。取同时包含全部 \(N_j\) 的有限 Galois 层 \(N\)。每个 \(U_{N_j}\) 包含 \(U_N\)，所以每个所选陪集以及它们的并，都是若干个 \(U_N\) 陪集的并。令 \(D=\pi_N(C)\)，便得等式。空集可直接取 \(D=\varnothing\)。
\end{solution}

\begin{exercise}\label{b1:ex:22-8}
设 \(L/K\) 为代数 Galois 扩张，\(G=\operatorname{Gal}(L/K)\) 取 Krull 拓扑，\(D\) 为离散群，不要求有限。证明每个连续同态 \(\rho:G\rightarrow D\) 的像有限，并且 \(\rho\) 经过某个有限 Galois 商 \(\operatorname{Gal}(N/K)\)。
\end{exercise}
\begin{solution}
由\cref{b1:thm:krull-compact}，\(G\) 紧，连续像 \(\rho(G)\) 也紧。离散空间的全部单点构成开覆盖，若其为无限集就没有有限子覆盖，所以 \(\rho(G)\) 必有限。将值域缩到这个有限离散群，\cref{b1:prop:finite-continuous-hom}给出一个有限 Galois 层 \(N\)，使 \(\rho\) 经过 \(G_N\)。还可直接从 \(\ker\rho=\rho^{-1}(1)\) 开并包含某个 \(U_N\) 得到同一结论。
\end{solution}

\begin{exercise}\label{b1:ex:22-9}
设 \(L/K\) 为代数 Galois 扩张，\(G=\operatorname{Gal}(L/K)\) 取 Krull 拓扑。若 \(H\leq G\) 为闭且指数有限的子群，证明 \(H\) 开。再取积拓扑群 \(V=\prod_{n\geq1}\mathbb F_2\) 与有限支集子群 \(W=\bigoplus_{n\geq1}\mathbb F_2\)，构造一个包含 \(W\) 的指数 \(2\) 真子群，并证明它不开。
\end{exercise}
\begin{solution}
左右平移为同胚，所以每个 \(H\) 陪集都闭。有限指数使 \(G\setminus H\) 只由有限多个这样的闭陪集组成，因此补集闭，\(H\) 开。对应于中间域的子群由\cref{b1:thm:infinite-galois-correspondence}保证闭，才可将此结论用于域论。

令 \(v=(1,1,\ldots)\in V\)，则 \(v\notin W\)。在 \(W\oplus\mathbb F_2v\) 上定义 \(\lambda_0(w+cv)=c\)，再把该子空间的一组基扩充为 \(V\) 的基，并规定 \(\lambda\) 在新增基向量上取零。这一步使用向量空间基的扩充，所得映射依赖所选的代数基，并不是由某个有限坐标规则给出的。它给出线性映射 \(\lambda:V\to\mathbb F_2\)。子群 \(H'=\ker\lambda\) 包含 \(W\)，不包含 \(v\)，且指数为 \(2\)。任意有限个坐标上的取值都可由 \(W\) 中的元素实现，故 \(W\) 稠密，\(H'\) 也稠密。若 \(H'\) 开，它的其他陪集也开，因而 \(H'\) 同时闭，与它稠密而非整个 \(V\) 矛盾。所以有限指数本身不能保证子群开。也可直接看出 \(\lambda\) 不连续：若连续，\(\lambda^{-1}(0)\) 应闭，包含稠密的 \(W\) 就必须等于 \(V\)，这与 \(\lambda(v)=1\) 矛盾。
\end{solution}

\begin{exercise}\label{b1:ex:22-10}
设 \(L/K\) 为代数 Galois 扩张，\(G=\operatorname{Gal}(L/K)\) 取 Krull 拓扑，\(K\subseteq E\subseteq L\) 且 \([E:K]<\infty\)。令 \(H=\operatorname{Gal}(L/E)\)。证明 \(C=\bigcap_{g\in G}gHg^{-1}\) 是开正规子群，其不动域是 \(E/K\) 在 \(L\) 内的正规闭包。
\end{exercise}
\begin{solution}
由\cref{b1:prop:finite-open-correspondence}，\(H\) 开且指数有限。群在有限集合 \(G/H\) 上作左乘，其核正是 \(C\)。这就是\cref{b1:prop:coset-action-core}中的正规核：它是包含在 \(H\) 中的最大 \(G\)-正规子群，因为若 \(J\trianglelefteq G\) 且 \(J\leq H\)，则 \(J=gJg^{-1}\leq gHg^{-1}\) 对所有 \(g\) 成立，故 \(J\leq C\)。只需选有限个陪集代表即可求这个交：若 \(g'=gh\)、\(h\in H\)，则 \(g'Hg'^{-1}=gHg^{-1}\)。因此 \(C\) 是有限个开子群的交，开且正规。

对任意 \(g\)，一个元素被 \(gHg^{-1}\) 固定当且仅当它属于 \(g(E)\)：把固定条件用 \(g^{-1}\) 共轭，应用\cref{b1:thm:infinite-galois-correspondence}即可。因此 \(C\) 恰为固定全部 \(g(E)\) 的自同构，\(L^C\) 就是这些域的合成域。它有限且在 \(K\) 上正规，因为 \(G\) 置换全部共轭域；任何包含 \(E\) 的正规中间扩张都包含这些共轭，所以这恰为最小的正规扩张，即正规闭包。子群一侧的“包含在 \(H\) 中且最大”，经包含反向的对应，成为域一侧的“包含 \(E\) 且最小”。
\end{solution}

\begin{exercise}\label{b1:ex:22-11}
令 \(G_{\mathbb F_q}=\operatorname{Gal}(\overline{\mathbb F}_q/\mathbb F_q)\cong\widehat{\vphantom{Z}\mathbb Z}\)，\(F(x)=x^q\) 为 Frobenius 自同构。取正整数 \(a,b\)，求 \(\langle F^a,F^b\rangle\) 的闭包，并证明
\[
a\widehat{\vphantom{Z}\mathbb Z}\cap b\widehat{\vphantom{Z}\mathbb Z}
=\operatorname{lcm}(a,b)\widehat{\vphantom{Z}\mathbb Z}.
\]
写出对应不动域的交与合成域。
\end{exercise}
\begin{solution}
在抽象循环子群 \(\langle F\rangle\) 中，Bézout 等式给出 \(\langle F^a,F^b\rangle=\langle F^d\rangle\)，其中 \(d=\gcd(a,b)\)。由\eqref{b1:eq:profinite-multiple-kernel}及其后的闭包计算，闭包为 \(d\widehat{\vphantom{Z}\mathbb Z}\)。

令 \(m=\operatorname{lcm}(a,b)\)。同一公式把左侧描述为模 \(a\)、模 \(b\) 同时为零的相容族。看其模 \(m\) 的坐标，一个整数剩余类同时被 \(a,b\) 整除当且仅当被 \(m\) 整除，所以左侧为模 \(m\) 的核，等于右侧。固定域反向对应遂给出
\[
\mathbb F_{q^a}\cap\mathbb F_{q^b}=\mathbb F_{q^d},
\qquad\mathbb F_{q^a}\mathbb F_{q^b}=\mathbb F_{q^m}.
\]
第一式对应生成子群的闭包，第二式对应闭子群的交。
\end{solution}

\begin{exercise}\label{b1:ex:22-12}
固定素数 \(\ell\)，令 \(E=\bigcup_{r\geq0}\mathbb F_{q^{\ell^r}}\subseteq\overline{\mathbb F}_q\)。证明 \(E/\mathbb F_q\) 是 Galois 扩张且 \(\operatorname{Gal}(E/\mathbb F_q)\cong\mathbb Z_\ell\)。在自然识别 \(\operatorname{Gal}(\overline{\mathbb F}_q/\mathbb F_q)\cong\widehat{\vphantom{Z}\mathbb Z}\) 下，求到 \(\operatorname{Gal}(E/\mathbb F_q)\) 的限制映射之核，并判断该核是否为零、是否开。
\end{exercise}
\begin{solution}
任一元素位于有限 Galois 层，其全部共轭都在该层中，故并域正规且可分。\cref{b1:prop:galois-compatible-family}把其 Galois 群识别为 \(\varprojlim_r\mathbb Z/\ell^r\mathbb Z\)，且有限坐标的基本开集一致，故为拓扑同构。由\cref{b1:thm:infinite-galois-quotient}，整个代数闭包上的限制满射。核为
\[
H=\bigcap_{r\geq1}\ell^r\widehat{\vphantom{Z}\mathbb Z}.
\]
每一项闭，所以 \(H\) 闭。这个交并不是零子群。对每个正整数 \(n\)，写 \(n=\ell^s m\)、\(\gcd(m,\ell)=1\)，用中国剩余定理选取唯一的剩余类
\[
e_n\equiv0\pmod{\ell^s},\qquad e_n\equiv1\pmod m.
\]
若 \(n\mid n'\)，则 \(e_{n'}\) 约化后仍满足这两个同余式，故 \((e_n)_n\) 是相容族。它在每个模 \(\ell^r\) 坐标为零，因而属于 \(H\)；但取任一不同于 \(\ell\) 的素数 \(p\)，其模 \(p\) 坐标为一，所以它非零。若 \(H\) 开，由\cref{b1:prop:finite-open-correspondence}，其固定域 \(E\) 应在 \(\mathbb F_q\) 上有限，与层次数 \(\ell^r\) 无界矛盾。因此该核不开。
\end{solution}

\begin{exercise}\label{b1:ex:22-13}
设 \(\ell\) 为素数，\(G_{\mathbb R}=\operatorname{Gal}(\mathbb C/\mathbb R)\) 取 Krull 拓扑。按基变换同构分类全部连续同态 \(G_{\mathbb R}\to\operatorname{GL}_2(\mathbb F_\ell)\)，其中目标群取离散拓扑。
\end{exercise}
\begin{solution}
由\cref{b1:prop:real-absolute-galois}，源群是离散的 \(C_2\)，所以每个表示都连续，由复共轭的像 \(A\) 决定，且 \(A^2=I\)。若 \(\ell\neq2\)，多项式 \(X^2-1\) 有不同根 \(1,-1\)，故空间分解为两个特征空间之和。按负特征空间的维数，恰有三个同构类，代表为 \(I\)、\(\operatorname{diag}(1,-1)\)、\(-I\)。

若 \(\ell=2\)，写 \(A=I+N\)，则 \(N^2=0\)。\(N=0\) 给出平凡表示；若 \(N\neq0\)，取 \(v\) 使 \(Nv\neq0\)，则 \(Nv,v\) 线性无关，否则 \(v\) 是 \(Nv\) 的倍数将推出 \(Nv=0\)。在这组基下 \(N\) 将第二个基向量 \(v\) 送到第一个基向量 \(Nv\)，并将 \(Nv\) 送到零，故 \(A\) 的代表为 \(\begin{psmallmatrix}1&1\\0&1\end{psmallmatrix}\)。因此特征 \(2\) 时恰有两个同构类。
\end{solution}

\begin{exercise}\label{b1:ex:22-14}
设 \(m,n\) 为正整数，\(\widehat{\vphantom{Z}\mathbb Z}\) 取逆极限拓扑，\(\mathbb Z/m\mathbb Z\) 取离散拓扑。求经过模 \(n\) 投影的连续同态 \(\widehat{\vphantom{Z}\mathbb Z}\to\mathbb Z/m\mathbb Z\) 的数目，其中有多少个是满射？再设 \(\ell\) 为素数，\(\rho:G_{\mathbb F_q}\to\operatorname{GL}_2(\mathbb F_\ell)\) 是把 Frobenius 自同构 \(F(x)=x^q\) 送到 \(A=\begin{psmallmatrix}1&1\\0&1\end{psmallmatrix}\) 的连续表示。对正整数 \(s\)，求其在 \(G_{\mathbb F_{q^s}}\) 上的限制，并求使该限制平凡的最小 \(s\)。所有绝对 Galois 群均在同一个 \(\overline{\mathbb F}_q\) 内取定。
\end{exercise}
\begin{solution}
由\cref{b1:prop:profinite-integer-finite-hom}，这样的连续同态由 \(1\) 的像 \(c\in\mathbb Z/m\mathbb Z\) 唯一决定。它经过模 \(n\) 投影，当且仅当 \(nc=0\)：这是 \(\mathbb Z/n\mathbb Z\) 的关系在目标中必须成立的条件。若 \(g=\gcd(m,n)\)，写 \(m=gm'\)、\(n=gn'\)，则 \(\gcd(m',n')=1\)。于是 \(m\mid nc\) 等价于 \(m'\mid n'c\)，再由互素性等价于 \(m'\mid c\)。模 \(m\) 的这些类为 \(0,m',\ldots,(g-1)m'\)，共有 \(g\) 个。满射还要求 \(c\) 的阶为 \(m\)，故当 \(m\mid n\) 时有 \(\varphi(m)\) 个满射，否则没有。

扩大有限基域以后，\(G_{\mathbb F_{q^s}}\) 是 \(G_{\mathbb F_q}\) 中的闭子群 \(s\widehat{\vphantom{Z}\mathbb Z}\)，它的拓扑生成元是 \(F^s\)。限制表示把这一生成元送到
\[
A^s=\begin{pmatrix}1&\bar s\\0&1\end{pmatrix},\qquad \bar s\in\mathbb F_\ell.
\]
连续性使这个取值唯一决定限制表示；具体地，它仍经过模 \(\ell\) 坐标，按右上角相加计算。若 \(\ell\mid s\)，则限制平凡；否则 \(A^s\) 的阶为 \(\ell\)，限制的像仍是由 \(A\) 生成的 \(\ell\) 阶群。因此使限制平凡的最小正整数是 \(s=\ell\)，相应基域为 \(\mathbb F_{q^\ell}\)。
\end{solution}

\begin{exercise}\label{b1:ex:22-15}
固定素数 \(\ell\)。对每个 \(r\geq1\)，记 \(a_{\ell^r}\) 为 \(a\in\widehat{\vphantom{Z}\mathbb Z}\) 的模 \(\ell^r\) 坐标，并定义
\[
\rho_r:\widehat{\vphantom{Z}\mathbb Z}\longrightarrow
\operatorname{GL}_2(\mathbb Z/\ell^r\mathbb Z),\qquad
a\longmapsto\begin{pmatrix}1&a_{\ell^r}\\0&1\end{pmatrix}.
\]
证明这些表示相容且连续，并求它们给出的 \(\mathbb Z_\ell\) 值矩阵表示的像与核。该表示是否经过单个有限商？
\end{exercise}
\begin{solution}
两个所列矩阵相乘时，右上角相加，其余位置不变，故 \(\rho_r\) 为同态；它经过有限坐标投影，所以连续。约化矩阵正好约化右上角，因而各层相容。根据正文在\subsecref{b1:subsec:2205:03}给出的逆极限构造，得到连续表示 \(\rho\)，其像为
\[
\Biggl\{\,\begin{pmatrix}1&b\\0&1\end{pmatrix}\ \Bigg|\ b\in\mathbb Z_\ell\,\Biggr\}.
\]
自然投影 \(\widehat{\vphantom{Z}\mathbb Z}\rightarrow\mathbb Z_\ell\) 在每一有限层都满射，因而由有限离散积的紧性及有限交性质可知它满射。故 \(\rho\) 的核为 \(\bigcap_r\ell^r\widehat{\vphantom{Z}\mathbb Z}\)。像是无限群，例如所有整数给出不同矩阵，因此不能经过单个有限商。每个有限坐标都经过有限层，并不意味着整族共有一个有限决定层。
\end{solution}

\begin{exercise}\label{b1:ex:22-16}
给 \(\overline{\mathbb F}_q\) 赋予离散拓扑，给 \(G_{\mathbb F_q}=\operatorname{Gal}(\overline{\mathbb F}_q/\mathbb F_q)\) 赋予 Krull 拓扑。证明 \(G_{\mathbb F_q}\) 对 \(\overline{\mathbb F}_q\) 的自然作用连续，每个元素的稳定子群都开，但整个作用的核不开。
\end{exercise}
\begin{solution}
若元素 \(a\) 在 \(\mathbb F_q\) 上次数为 \(d\)，则 \(\mathbb F_q(a)=\mathbb F_{q^d}\)，其稳定子群是模 \(d\) 投影的核 \(d\widehat{\vphantom{Z}\mathbb Z}\)，由\eqref{b1:eq:profinite-multiple-kernel}为开子群。对任意 \((g,a)\)，集合 \(g\operatorname{Stab}(a)\times\{a\}\) 是一个开邻域，整个邻域上的作用值均为 \(g(a)\)，因而作用连续。

整体核必须固定全部代数元，所以只有恒等自同构。若它开，则由\cref{b1:prop:finite-open-correspondence}，其固定域 \(\overline{\mathbb F}_q\) 应是有限扩张，矛盾。这里对每个元素分别找到开稳定子群，不能换成“存在一个开子群同时固定所有元素”。
\end{solution}

\begin{exercise}\label{b1:ex:22-17}
设 \(L/K\) 为代数 Galois 扩张，\(G=\operatorname{Gal}(L/K)\) 取 Krull 拓扑，\(H,J\leq G\) 为闭子群。若 \(N/K\) 是 \(L/K\) 的有限 Galois 中间扩张，\(\pi_N:G\to\operatorname{Gal}(N/K)\) 为限制映射，是否必有 \(\pi_N(H\cap J)=\pi_N(H)\cap\pi_N(J)\)？用一个有限 Galois 扩张给出反例。
\end{exercise}
\begin{solution}
取 \(L=\mathbb Q(\sqrt2,\sqrt3)\)。\cref{b1:prop:biquadratic-example}给出 \([L:\mathbb Q]=4\)；它又是可分多项式 \((X^2-2)(X^2-3)\) 的分裂域，所以\cref{b1:prop:galois-automorphism-count}给出 \(|G|=4\)。每个自同构由两个平方根的像决定，像分别只能为原平方根或其负数，故取符号给出单射 \(G\rightarrow C_2^2\)。两边同有四个元素，因而是同构。以这两个符号为坐标，令 \(H=\langle(1,1)\rangle\)、\(J=\langle(1,0)\rangle\)，并令 \(\pi\) 为限制到 \(\mathbb Q(\sqrt2)\)，即取第一坐标。

这时 \(H\cap J=0\)，故 \(\pi(H\cap J)=0\)；但两个子群各自都满射到第一个 \(C_2\)，所以 \(\pi(H)\cap\pi(J)=C_2\)。所有子群在有限离散群中都闭，反例仍成立。\cref{b1:prop:closure-finite-quotients}使用的是全部有限层的相容信息，并未断言单个投影保持子群的交。
\end{solution}

\begin{exercise}\label{b1:ex:22-18}
将 \(\mathbb Z\) 视为它在 \(\widehat{\vphantom{Z}\mathbb Z}\) 中的自然像，令 \(Q=\widehat{\vphantom{Z}\mathbb Z}/\mathbb Z\) 带商拓扑。证明任何从 \(Q\) 到 Hausdorff 空间的连续映射都是常值映射，并解释为何这不说明抽象群 \(Q\) 平凡。
\end{exercise}
\begin{solution}
先直接核对商拓扑。非空开集的原像 \(V\subseteq\widehat{\vphantom{Z}\mathbb Z}\) 是非空开集且满足 \(V+\mathbb Z=V\)。给任意 \(g\)，开集 \(g-V\) 与稠密的 \(\mathbb Z\) 相交，故 \(g\in V+\mathbb Z=V\)。所以商中只有空集与全空间两个开集。

若连续映射取两个不同值，Hausdorff 性给出分别包含这两个值的不相交开邻域，其原像是两个非空不相交开集，与上述描述矛盾，故映射常值。正文在\subsecref{b1:subsec:2203:02}构造的相容族，在所有模 \(2^r\) 坐标为零、在模 \(3\) 坐标为一，不能来自整数。因此 \(\mathbb Z\) 是真子群，抽象商非平凡。失去的是用 Hausdorff 拓扑区分这些陪集的能力。
\end{solution}

\begin{exercise}\label{b1:ex:22-19}
设 \(L/K\) 为代数 Galois 扩张，\(G=\operatorname{Gal}(L/K)\) 取 Krull 拓扑，\(H\) 为 \(G\) 的闭正规子群。对 \(L/K\) 的每个有限 Galois 中间域 \(N/K\)，记 \(G_N=\operatorname{Gal}(N/K)\)，并记限制映射为 \(\pi_N:G\to G_N\)。证明有限群 \(G_N/\pi_N(H)\) 构成逆系统，并证明自然映射
\[
G/H\longrightarrow\varprojlim_N G_N/\pi_N(H)
\]
是拓扑群同构。
\end{exercise}
\begin{solution}
\(\pi_N(H)\) 正规；限制映射将 \(\pi_M(H)\) 满射到 \(\pi_N(H)\)，所以诱导的商映射良定义且相容。先定义连续同态
\[
\Phi:G\longrightarrow\varprojlim_NG_N/\pi_N(H),
\qquad g\longmapsto\bigl(\pi_N(g)\pi_N(H)\bigr)_N.
\]
各坐标相容，故 \(\Phi\) 确实取值于逆极限；每个坐标连续，故 \(\Phi\) 连续。由\cref{b1:prop:closure-finite-quotients}，其核为
\[
\ker\Phi=\bigcap_N\pi_N^{-1}\bigl(\pi_N(H)\bigr)=\overline H=H.
\]
因此 \(\Phi\) 诱导单同态 \(\overline\Phi:G/H\to\varprojlim_NG_N/\pi_N(H)\)。由商拓扑的定义，\(\overline\Phi\) 连续。

给定右侧的相容陪集族，对每个 \(N\) 令 \(C_N\subseteq G\) 为这一陪集在 \(G\) 中的原像。限制满射使 \(C_N\) 非空，它又开且闭。任取有限个 \(N\)，选共同上层 \(M\)；相容性保证 \(C_M\) 包含于它们的交。所以 \((C_N)_N\) 有有限交性质，由\cref{b1:thm:krull-compact}总交非空。这证明满射。

\(G\) 紧，故其商空间 \(G/H\) 也紧；目标是有限离散群直积的子空间，具有 Hausdorff 性。\cref{b1:prop:compact-maps}说明连续双射 \(\overline\Phi\) 为同胚，得到所需拓扑同构。闭性在核由 \(\overline H\) 变为 \(H\) 的步骤不可省略。
\end{solution}

\begin{exercise}\label{b1:ex:22-20}
设 \(T\) 是有限零一字组成的集合，对取前缀封闭，且每个长度都有 \(T\) 中的字。证明存在无限零一序列，其每个有限前缀都属于 \(T\)。先用有限离散积的紧性证明，再给出按次序选 \(0\) 或 \(1\) 的直接构造，说明此特殊情形无需为每一步作无规则的选择。
\end{exercise}
\begin{solution}
在 \(\{0,1\}^{\mathbb N}\) 中，令 \(C_r\) 为前 \(r\) 位组成的字属于 \(T\) 的序列。它是有限多个基本开闭集的并，故闭且非空；非空性可将任一长度为 \(r\) 的字后补零。前缀封闭使 \(C_{r+1}\subseteq C_r\)，所以具有有限交性质。\cref{b1:lem:finite-discrete-product-compact}给出非空总交。

直接构造时，如果一个字在 \(T\) 中具有任意长的延长字，就称它可以无限延长。空字有此性质。若一个字有此性质，它的两个一步延长至少一个仍有此性质；否则两个分支的长度各有上界，取较大上界便产生矛盾。规定能选 \(0\) 分支时就选 \(0\)，否则选 \(1\)，递归产生所需无限序列。这个规则在每一步唯一确定所选的分支。
\end{solution}
